%!TEX root = forallxyyc.tex
\part{真值函项逻辑的自然演绎}
\label{ch.NDTFL}
\addtocontents{toc}{\protect\mbox{}\protect\hrulefill\par}

\chapter{自然演绎的基本思想}\label{s:NDVeryIdea}

早在 \cref{s:Valid} 中，我们就说过，一个论证是有效的 \ifeff{} 不存在所有前提都为真而结论为假的情形。

在真值函项逻辑中，这引导我们发展出真值表。一个完整真值表的每一行都对应一个赋值。因此，面对一个真值函项逻辑论证时，要判断是否存在某个赋值使得所有前提为真而结论为假，我们有一个非常直接的方法：只需完整地构造真值表即可。

然而，真值表可能无法提供多少\emph{洞见}。考虑真值函项逻辑中的两个论证：
\begin{align*}
P \eor Q, \enot P & \therefore Q\\
P \eif Q, P & \therefore Q
\end{align*}
显然，这些都是有效论证。你可以通过构造四行真值表来确认这一点，但我们或许会说，它们利用了不同的\emph{推理形式}。若能跟踪这些不同的推理形式，会很有帮助。

\emph{自然演绎系统}的目标之一，就在于以能够让我们理解论证所涉推理过程的方式，来表明特定论证是有效的。我们从非常基本的推理规则开始。这些规则可以组合起来，以构建更复杂的论证。事实上，我们希望仅用一小套初始的推理规则，就能刻画所有有效论证。

\emph{这是思考论证的一种非常不同的方式。}

使用真值表时，我们直接考虑使语句为真或为假的各种方式。使用自然演绎系统时，我们则依据我们认定为良好的规则来操作语句。后者预示着能让我们更深入地理解——或者至少是提供一种不同的视角来理解——论证是如何运作的。

转向自然演绎的动机，可能不仅源于寻求洞见，也可能是出于\emph{必要性}。考虑：
\ifHTMLtarget


\begin{earg}
	\item $A_1 \eif C_1$
	\item[\texttherefore] $(A_1 \eand (A_2 \eand (A_3 \eand (A_4 \eand A_5)))) \eif
	((((C_1 \eor C_2) \eor C_3) \eor C_4) \eor C_5)$
\end{earg}


\else
\begin{align*}
	& A_1 \eif C_1 \\
	\therefore\, & (A_1 \eand (A_2 \eand (A_3 \eand (A_4 \eand A_5)))) \eif {}\\
	& \qquad ((((C_1 \eor C_2) \eor C_3) \eor C_4) \eor C_5)
\end{align*}
\fi
要检验这个论证是否有效，你可能需要构造一个 1{,}024 行的真值表。如果你正确地完成了它，会发现没有一行是所有前提为真而结论为假的。因此你会知道这个论证是有效的。（但如前所述，在某种意义上你并不知道\emph{为什么}这个论证有效。）但现在考虑：
\ifHTMLtarget


\begin{earg}
\item $A_1 \eif C_1$
\item[\texttherefore] $(A_1 \eand (A_2 \eand (A_3 \eand (A_4 \eand (A_5 \eand
 (A_6 \eand (A_7 \eand (A_8 \eand (A_9 \eand A_{10}))))))))) \eif
 (((((((((C_1 \eor C_2) \eor C_3) \eor C_4) \eor C_5) \eor  C_6) \eor
 C_7) \eor C_8) \eor C_9) \eor C_{10})$
\end{earg}


\else
\begin{align*}
& A_1 \eif C_1 \\
\therefore\, & \small (A_1 \eand (A_2 \eand (A_3 \eand (A_4 \eand (A_5 \eand {}\\
&\qquad (A_6 \eand (A_7 \eand (A_8 \eand (A_9 \eand A_{10}))))))))) \eif {}\\
&(((((((((C_1 \eor C_2) \eor C_3) \eor C_4) \eor C_5) \eor {}\\
& \qquad C_6) \eor C_7) \eor C_8) \eor C_9) \eor C_{10})
\end{align*}
\fi
这个论证（你或许能看出来）也是有效的——但要检验它需要一个包含 $2^{20} = 1{,}048{,}576$ 行的真值表。原则上，我们可以让机器机械地检查真值表，完成后再报告结果。但实际上，如果使用真值表，真值函项逻辑中的复杂论证会变得\emph{难以处理}。

然而，当我们进入一阶逻辑（从 \cref{s:FOLBuildingBlocks} 开始）时，问题会急剧恶化。一阶逻辑没有类似真值表的测试方法。要评估一个论证是否有效，我们必须对\emph{所有}解释进行推理，但正如我们将要看到的，可能的解释有无穷多个。我们甚至在原则上也无法让一台机器遍历所有无穷多个可能的解释并在完成后报告结果：它\emph{永远}无法完成。我们或者需要找到一种更高效的方法来对所有解释进行推理，或者必须寻找其他途径。

确实存在一些系统，能够将关于所有可能解释的推理方法规范化。它们由 Evert Beth 和 Jaakko Hintikka 在 20 世纪 50 年代提出，但我们不会遵循这条路径。相反，我们将着眼于自然演绎。

我们将不再直接对所有赋值（在真值函项逻辑的情况下）进行推理，而是试图选取一些基本的推理规则。其中一些规则将支配语句联结词的行为。另一些则支配构成一阶逻辑特色的量词和等词的行为。由此产生的规则系统将为我们提供一种思考论证有效性的新方式。
现代自然演绎的发展源于 Gerhard Gentzen 和 Stanis\l{}aw Ja\'{s}kowski 同时且不相关的论文（1934）。不过，我们将要考察的自然演绎系统主要基于 Frederic Fitch 的工作（首次发表于 1952 年）。

\chapter{真值函项逻辑的基本规则}\label{s:BasicTFL}
我们将建立一个 \define{自然演绎} 系统。对于每个联结词，都有一套 \define{引入规则}（允许我们证明以该联结词为主逻辑算子的语句）和 \define{消除规则}（允许我们从以该联结词为主逻辑算子的语句出发证明某些东西）。

\section{形式推导的概念}
一个 \emph{形式推导} 或 \emph{形式证明} 是一个语句序列，其中某些语句被标记为初始假设（或前提）。形式推导的最后一行是结论。（此后，我们将简称为“推导”或“证明”，但请注意，也存在\emph{非形式推导}。）

举例来说，考虑：
	$$\enot (A \eor B) \therefore \enot A \eand \enot B$$
我们首先写下前提来开始推导：

\begin{fitchproof}
	\hypo{a1}{\enot (A \eor B)}\PR
\end{fitchproof}

注意，我们为这个前提编了号，因为后面可能需要引用它。实际上，推导的每一行都有编号，以便我们能够回指。

还要注意，我们在前提下面划了一条线。线上方所写的都是\emph{假设}。线下方所写的，要么是从假设推出的结果，要么是某个新的假设。我们期望最终得到结论“$\enot A \eand \enot B$”；因此我们希望推导最终以下面这样一行结束：

\begin{fitchproof}
	\have[n]{con}{\enot A \eand \enot B}
\end{fitchproof}

其中 $n$ 是某个行号。具体在哪一行结束并不重要，但我们显然希望推导越短越好。

类似地，假设我们想考虑：
$$A\eor B, \enot (A\eand C), \enot (B \eand \enot D) \therefore \enot C\eor D$$
这个论证有三个前提，所以一开始我们把它们全部写下来、编号，并在下方划线：

\begin{fitchproof}
	\hypo{a1}{A \eor B}\PR
	\hypo{a2}{\enot (A\eand C)}\PR
	\hypo{a3}{\enot (B \eand \enot D)}\PR
\end{fitchproof}

我们希望最终得到某一行：

\begin{fitchproof}
	\have[n]{con}{\enot C \eor D}
\end{fitchproof}

剩下的工作，就是要解释在从前提到结论的过程中我们可以使用的每一条规则。这些规则依照逻辑联结词来分类。

\section{复述规则}
第一条规则如此显而易见，以至于我们竟然要专门说明它，这都令人惊讶。

如果你已经在一个推导中得到了某个语句，\emph{复述规则}允许你在新的一行上重复它。例如：

\begin{fitchproof}
	\have[4]{a1}{A \eand B}
	\ellipsesline
	\have[10]{a2}{A \eand B} \by{R}{a1}
\end{fitchproof}

这表明我们在第 $4$ 行写出了“$A \eand B$”。现在，在某个后续行——例如第 $10$ 行——我们决定要重复它。因此我们再次写出它。我们还需添加一个引用来证成我们写下的内容。这里，我们写“R”表示我们使用了复述规则，写“$4$”表示我们将其应用于第 $4$ 行。

以下是该规则的一般表述：
\factoidbox{


\begin{fitchproof}
	\have[m]{a}{\metav{A}}
	\have[\ ]{c}{\metav{A}} \by{R}{a}
\end{fitchproof}

}
要点在于，如果某个语句 $\metav{A}$ 出现在某一行，那么我们可以在后续行重复 $\metav{A}$。我们证明的每一行都必须有一个证成，此处证成是“R $m$”。这表示：复述，应用于第 $m$ 行。

有两点需要强调。第一，“$\metav{A}$”不是真值函项逻辑的语句。它是元语言中的一个符号，用于代表任意一个真值函项逻辑语句（参见 \cref{s:UseMention}）。
第二，类似地，“$m$”也不是会出现在证明中的符号。它是元语言中的一个符号，用于指代证明中的任意一个行号。在实际的证明中，行号是“$1$”、“$2$”、“$3$”等等。但在定义规则时，我们使用像“$m$”这样的变量来强调这条规则可以在任意点应用。

\section{合取}
假设我们想证明 Ludwig 既是反动派又是自由意志主义者。一种显而易见的做法如下：首先证明 Ludwig 是反动派；然后证明 Ludwig 是自由意志主义者；接着将这两个证明结合起来，得到合取式。

我们的自然演绎系统将直接把握这一思路。在上述例子中，我们可以采用如下符号化约定：

\begin{ekey}
		\item[R] Ludwig 是反动派
		\item[L] Ludwig 是自由意志主义者
	\end{ekey}

假设我们正在进行一个推导，并且已经在第 $8$ 行得到了“$R$”，在第 $15$ 行得到了“$L$”。那么，在任意后续行，我们可以这样得到“$R \eand L$”：

\begin{fitchproof}
	\have[8]{a}{R}
	\have[15]{b}{L}
	\have[\ ]{c}{R \eand L} \ai{a, b}
\end{fitchproof}

注意，我们证明中的每一行要么是一个假设，要么必须由某条规则证成。这里我们引用“$\eand$I $8$, $15$”表明该行是通过将合取引入规则（$\eand$I）应用于第 $8$ 行和第 $15$ 行而得到的。我们同样可以得到：


\begin{fitchproof}
	\have[8]{a}{R}
	\have[15]{b}{L}
	\have[\ ]{c}{L \eand R} \ai{b, a}
\end{fitchproof}


其中的引用顺序是颠倒的，以反映合取支的顺序。更一般地说，以下是我们的合取引入规则：
\factoidbox{


\begin{fitchproof}
	\have[m]{a}{\metav{A}}
	\have[n]{b}{\metav{B}}
	\have[\ ]{c}{\metav{A}\eand\metav{B}} \ai{a, b}
\end{fitchproof}

}

再次提醒，这个规则的陈述是\emph{图式性}的。它本身不是一个证明。“$\metav{A}$”和“$\metav{B}$”是代表真值函项逻辑中任意语句的元变元。类似地，“$m$”和“$n$”不是会在任何实际证明中出现的数字，而是我们用来指称任意证明中任意行号的符号。我们使用这类元变元来强调该规则可以在任何地方应用。规则只要求两个合取支在证明中更早的某个地方可为我们所用。它们可以彼此分离，也可以以任何顺序出现。

这条规则被称为“合取\emph{引入}”规则，因为它将符号“$\eand$”引入我们的证明，而此前该符号可能并不存在。相应地，我们有一条\emph{消去}该符号的规则。假设你已证明路德维希既是反动派又是自由意志主义者。那么你可以合理得出结论：路德维希是反动派。同样，你也可以得出结论：路德维希是自由意志主义者。结合起来，我们便得到合取消去规则：
\factoidbox{


\begin{fitchproof}
	\have[m]{ab}{\metav{A}\eand\metav{B}}
	\have[\ ]{a}{\metav{A}} \ae{ab}
\end{fitchproof}

}
同样地：
\factoidbox{


\begin{fitchproof}
	\have[m]{ab}{\metav{A}\eand\metav{B}}
	\have[\ ]{b}{\metav{B}} \ae{ab}
\end{fitchproof}

}
其要点很简单：当你在证明的某一行有一个合取式时，你可以使用 $\eand$E 规则得到其中任何一个合取支。有一点值得强调：你只能在合取式的主逻辑算子是“$\eand$”时应用此规则。因此，你不能仅从“$C \eor (D \eand E)$”就推出“$D$”！

即便只有这两条规则，我们也已能初步领略形式证明系统的威力。考虑如下论证：


\begin{earg}
\item[] $[(A\eor B)\eif(C\eor D)] \eand [(E \eor F) \eif (G\eor H)]$
\item[\texttherefore] $[(E \eor F) \eif (G\eor H)] \eand [(A\eor B)\eif(C\eor D)]$
\end{earg}


此论证前提和结论的主逻辑算子都是“$\eand$”。为构造证明，我们首先写下作为假设的前提，并在其下方画一条横线：此线之后的所有内容都必须（通过反复应用推理规则）从我们的假设中推出。因此证明的开头如下：


\begin{fitchproof}
	\hypo{ab}{{[}(A\eor B)\eif(C\eor D){]} \eand {[}(E \eor F) \eif (G\eor H){]}}\PR
\end{fitchproof}


从前提出发，我们可运用 $\eand$E 得到每一个合取支。此时证明如下：


\begin{fitchproof}
	\hypo{ab}{{[}(A\eor B)\eif(C\eor D){]} \eand {[}(E \eor F) \eif (G\eor H){]}}\PR
	\have{a}{{[}(A\eor B)\eif(C\eor D){]}} \ae{ab}
	\have{b}{{[}(E \eor F) \eif (G\eor H){]}} \ae{ab}
\end{fitchproof}


于是，通过将 $\eand$I 规则应用于第 3 行和第 2 行（按此顺序），我们便得到所需结论。完整的证明如下：


\begin{fitchproof}
	\hypo{ab}{{[}(A\eor B)\eif(C\eor D){]} \eand {[}(E \eor F) \eif (G\eor H){]}}\PR

	\have{a}{{[}(A\eor B)\eif(C\eor D){]}} \ae{ab}
	\have{b}{{[}(E \eor F) \eif (G\eor H){]}} \ae{ab}
	\have{ba}{{[}(E \eor F) \eif (G\eor H){]} \eand {[}(A\eor B)\eif(C\eor D){]}} \ai{b,a}
\end{fitchproof}


这是一个非常简单的证明，但它展示了如何将证明规则链接起来构成更长的证明。顺便指出，若用真值表检验此论证需要 256 行；而我们的形式证明仅需四行。

值得再举一例。回顾 \cref{s:MoreBracketingConventions}，我们曾指出以下论证是有效的：
	$$A \eand (B \eand C) \therefore (A \eand B) \eand C$$
为给出相应证明，我们首先写下：


\begin{fitchproof}
	\hypo{ab}{A \eand (B \eand C)}\PR
\end{fitchproof}


从前提出发，通过两次应用 $\eand$E，可得到每个合取支。接着可再应用两次 $\eand$E，证明如下：


\begin{fitchproof}
	\hypo{ab}{A \eand (B \eand C)}\PR
	\have{a}{A} \ae{ab}
	\have{bc}{B \eand C} \ae{ab}
	\have{b}{B} \ae{bc}
	\have{c}{C} \ae{bc}
\end{fitchproof}


现在，我们可以按所需顺序顺利地重新引入合取，从而得到最终证明：


\begin{fitchproof}
	\hypo{abc}{A \eand (B \eand C)}\PR
	\have{a}{A} \ae{abc}
	\have{bc}{B \eand C} \ae{abc}
	\have{b}{B} \ae{bc}
	\have{c}{C} \ae{bc}
	\have{ab}{A \eand B}\ai{a, b}
	\have{con}{(A \eand B) \eand C}\ai{ab, c}
\end{fitchproof}


回想一下，真值函项逻辑中语句的正式定义只允许包含两个合取支的合取式。上述证明似乎暗示我们可以在所有证明中省略内部括号。然而，这并非标准做法，我们也不这样做。相反，我们将保持更为简洁的括号约定（尽管为了可读性，我们仍允许省略最外层括号）。

最后再举一例。运用 $\eand$I 规则时，并不要求将其应用于不同的语句。因此，若我们愿意，可以从“$A$”形式地证明“$A \eand A$”如下：


\begin{fitchproof}
	\hypo{a}{A}\PR
	\have{aa}{A \eand A}\ai{a, a}
\end{fitchproof}


简单而有效。

\section{条件句}
考虑以下论证：


\begin{earg}
		\item[] 如果简聪明，那么她敏捷。
		\item[] 简聪明。
		\item[\texttherefore] 简敏捷。
\end{earg}


此论证显然是有效的，它直接提示了条件消去规则（$\eif$E）：
\factoidbox{


\begin{fitchproof}
	\have[m]{ab}{\metav{A}\eif\metav{B}}
	\have[n]{a}{\metav{A}}
	\have[\ ]{b}{\metav{B}} \ce{ab,a}
\end{fitchproof}

}
此规则有时亦称为\emph{肯定前件式（modus ponens）}。同样，这是一条消去规则，因为它允许我们从包含“$\eif$”的语句出发，得到一个可能不包含“$\eif$”的语句。注意，条件式 $\metav{A}\eif\metav{B}$ 与其前件 $\metav{A}$ 在证明中可以彼此分离，且可以任何顺序出现。惯例是先引用包含条件式 $\metav{A}\eif\metav{B}$ 的行号。

条件引入规则的动机也相当直接。下述论证应是有效的：


\begin{earg}
		\item 路德维希是反动派。
		\item[\texttherefore] 如果路德维希是自由意志主义者，那么路德维希既是反动派\emph{又}是自由意志主义者。
	\end{earg}


若有人质疑其有效性，我们或可如下解释以说服对方：


\begin{quote}
		假设路德维希是反动派。现在，\emph{额外再}假设路德维希是自由意志主义者。那么根据刚才讨论的合取引入规则，路德维希既是反动派又是自由意志主义者。当然，这依赖于“路德维希是自由意志主义者”这一假设。但这恰恰意味着：如果路德维希是自由意志主义者，那么他既是反动派又是自由意志主义者。
	\end{quote}


转换为自然演绎格式，我们刚才使用的推理模式如下。我们从单一前提“路德维希是反动派”开始：


\begin{fitchproof}
		\hypo{r}{R}\PR
	\end{fitchproof}


接下来，为了论证，我们做一个\emph{额外的}假设（“路德维希是自由意志主义者”）。为了表明我们不再仅仅处理原始前提（“$R$”），而是处理一个额外假设，我们继续证明如下：


\begin{fitchproof}
		\hypo{r}{R}\PR
		\open
			\hypo{l}{L}\AS
	\end{fitchproof}


注意，在第 $2$ 行我们\emph{并非}声称从第 $1$ 行证明了“$L$”，因此没有为该额外假设标注任何证立依据。但我们需要标明这是一个额外假设：我们在其下方画一横线（表示它是一个假设），并用一条额外的垂直线使其缩进（表示它是额外的）。

有了这个额外假设，我们便可使用 $\eand$I。于是可以继续证明：


\begin{fitchproof}
		\hypo{r}{R}\PR
		\open
			\hypo{l}{L}\AS
			\have{rl}{R \eand L}\ai{r, l}
%			\close
%		\have{con}{L \eif (R \eand L)}\ci{l-rl}
	\end{fitchproof}


这样，我们已表明在额外假设“$L$”下，可以得到“$R \eand L$”。因此我们可以得出结论：若“$L$”成立，则“$R \eand L$”也成立。或者更简洁地说，我们可以得到“$L \eif (R \eand L)$”：


\begin{fitchproof}
		\hypo{r}{R}\PR
		\open
			\hypo{l}{L}\AS
			\have{rl}{R \eand L}\ai{r, l}
			\close
		\have{con}{L \eif (R \eand L)}\ci{l-rl}
	\end{fitchproof}


注意，左侧已恢复为单条垂直线。我们已\emph{解除}了额外假设“$L$”，因为条件句本身仅从原始前提“$R$”即可得出。

这里运作的一般模式如下。我们先做出一个附加假设 $\metav{A}$；然后从那个附加假设出发，证明 $\metav{B}$。这样一来，我们就知道了下面这一点：如果 $\metav{A}$ 为真，那么 $\metav{B}$ 为真。这被总结为条件引入规则：
\factoidbox{


\begin{fitchproof}
		\open
			\hypo[i]{a}{\metav{A}}\AS
			\have[j]{b}{\metav{B}}
		\close
		\have[\ ]{ab}{\metav{A}\eif\metav{B}}\ci{a-b}
\end{fitchproof}

}
在 $i$ 行和 $j$ 行之间，你可以使用任意多行。

再举一个 $\eif$I 规则的实际例子会很有帮助。假设我们想考虑以下论证：
	$$P \eif Q, Q \eif R \therefore P \eif R$$
我们首先列出\emph{全部}两个前提。接着，因为我们的目标是得到一个条件式（即 “$P \eif R$”），所以我们额外假设那个条件式的前件。于是，我们的主证明这样开始：


\begin{fitchproof}
	\hypo{pq}{P \eif Q}\PR
	\hypo{qr}{Q \eif R}\PR
	\open
		\hypo{p}{P}\AS
	\close
\end{fitchproof}


注意，我们通过将 “$P$” 作为一个附加假设，使之可用。据此，我们可以对第一个前提使用 {\eif}E 规则，得到 “$Q$”。接着，我们可以对第二个前提使用 {\eif}E 规则。所以，通过假设 “$P$” 我们得以证明 “$R$”，于是我们应用 {\eif}I 规则——解除 “$P$” 这一假设——并完成证明。把所有这些放在一起，我们得到：
\label{HSproof}


\begin{fitchproof}
	\hypo{pq}{P \eif Q}\PR
	\hypo{qr}{Q \eif R}\PR
	\open
		\hypo{p}{P}\AS
		\have{q}{Q}\ce{pq,p}
		\have{r}{R}\ce{qr,q}
	\close
	\have{pr}{P \eif R}\ci{p-r}
\end{fitchproof}

\section{附加假设与子证明}
$\eif$I 规则运用了做出附加假设的思想。这些假设需要小心处理。考虑这个证明：


\begin{fitchproof}
	\hypo{a}{A}\PR
	\open
		\hypo{b1}{B}\AS
		\have{b2}{B} \by{R}{b1}
	\close
	\have{con}{B \eif B}\ci{b1-b2}
\end{fitchproof}

这完全符合我们已经制定的规则，而且不应该显得特别奇怪。既然 “$B \eif B$” 是一个重言式，证明它自然不需要任何特定的前提。

但假设我们现在试图如下继续这个证明：

\begin{fitchproof}
	\hypo{a}{A}\PR
	\open
		\hypo{b1}{B}\AS
		\have{b2}{B} \by{R}{b1}
	\close
	\have{con}{B \eif B}\ci{b1-b2}
\ifHTMLtarget
	\have{b}{B} \by{不当的尝试以应用$\eif$E}{con, b2}
\else
	\have{b}{B} \by{不当的尝试}{}
	\have [\ ]{x}{} \by{以应用$\eif$E}{con, b2}
\fi
\end{fitchproof}

如果我们被允许这样做，那将是一场灾难。那将允许我们由任一语句字母推导出任何其他语句字母。然而，如果你告诉我安妮很快（用 “$A$” 符号化），我们不应该能够得出结论说布狄卡女王身高二十英尺（用 “$B$” 符号化）！我们必须禁止这样做，但我们如何实施这一禁止呢？

我们可以把引入额外假设的过程描述为进行一次\emph{子证明}：即主证明内部的一个辅助证明。开始子证明时，我们会另画一条竖线，表示我们不再处于主证明中。然后，我们写下该子证明所基于的假设。实质上，我们可以把子证明看作是在提出这样一个问题：\emph{如果我们也做出这个额外假设，那么我们可以证明什么？}

在子证明内部进行推导时，我们可以引用引入该子证明时所引入的额外假设，也可以引用从我们最初的那些假设得出的任何结果。（毕竟，那些最初的假设仍然有效。）然而，在某个时刻，我们会希望停止使用这个额外假设：我们希望从子证明返回到主证明。为了表明我们已经返回到主证明，子证明的竖线到此终止。此时，我们说该子证明是\define{闭合的（closed）}。当我们闭合一个子证明时，我们就解除了那个额外假设，因此再依赖任何依赖于该额外假设的东西都是不合法的。于是我们规定：
\factoidbox{在应用规则时，若要引用某一行：

\begin{compactlist}
\item 该行必须出现在应用规则的那一行之前，但是
\item 不能出现在一个在应用规则的那一行之前就已经闭合的子证明之内。
\end{compactlist}

}
这条规定排除了上面那个灾难性的证明尝试。在第5行应用规则$\eif$E要求我们引用证明中更早出现的两个独立行。其中一行（即第3行）出现在一个子证明（第2-3行）内部。到我们必须引用第3行的第5行时，那个子证明已经闭合了。这是不合法的：我们不能在第5行引用第3行。

闭合一个子证明被称为\define{解除}该子证明的假设。
因此我们可以这样概括要点：\emph{你不可以回头引用任何依赖于已解除假设所得到的东西}。

那么，子证明让我们能够思考，如果作出额外假设，我们能证明什么。从这里得出的要点并不令人惊讶——在证明过程中，我们必须非常仔细地记录在任一给定时刻我们作了哪些假设。我们的证明系统以非常图示化的方式做到这一点。（确实，这正是我们选择使用\emph{这个}证明系统的原因。）

一旦我们开始思考通过增加假设能证明什么，就没有什么能阻止我们提出这个问题：如果我们作出\emph{更多}假设，我们能证明什么？这可能促使我们在一个子证明内部再引入一个子证明。以下是一个仅使用我们目前考虑过的规则的例子：


\begin{fitchproof}
\hypo{a}{A}\PR
\open
	\hypo{b}{B}\AS
	\open
		\hypo{c}{C}\AS
		\have{ab}{A \eand B}\ai{a,b}
	\close
	\have{cab}{C \eif (A \eand B)}\ci{c-ab}
\close
\have{bcab}{B \eif (C \eif (A \eand B))}\ci{b-cab}
\end{fitchproof}


注意第4行的引用回溯到了初始假设（第1行）和一个子证明的假设（第2行）。这完全符合规则，因为至该行（即第4行）为止，两个假设都尚未被解除。

不过，我们再次需要仔细记录在任一给定时刻我们所作的假设。假设我们试图如下继续这个证明：


\begin{fitchproof}
\hypo{a}{A}\PR
\open
	\hypo{b}{B}\AS
	\open
		\hypo{c}{C}\AS
		\have{ab}{A \eand B}\ai{a,b}
	\close
	\have{cab}{C \eif (A \eand B)}\ci{c-ab}
\close
\have{bcab}{B \eif(C \eif (A \eand B))}\ci{b-cab}
\ifHTMLtarget
	\have{bcab}{C \eif (A \eand B)}\by{不当的尝试以应用 $\eif$I}{c-ab}
\else
	\have{bcab}{C \eif (A \eand B)}\by{不当的尝试}{}
	\have [\ ]{x}{} \by{以应用 $\eif$I}{c-ab}
\fi
\end{fitchproof}


那就太糟糕了。如果我们告诉你 Anne 是聪明的，你不应该能推断出：如果 Cath 是聪明的（用 “$C$” 符号化），那么\emph{既} Anne 是聪明的\emph{又} Queen Boudica 身高 20 英尺！但如果这种证明是允许的，那它恰恰就会暗示这样的结论。

核心问题在于，以假设 “$C$” 开始的子证明关键依赖于我们在第 2 行已经假设了 “$B$”。到第 6 行，我们已经\emph{解除}了假设 “$B$”：我们已经不再追问如果也假设 “$B$” 我们能证明什么了。因此，试图（在第 7 行）利用以假设 “$C$” 开始的子证明，这纯属作弊。于是我们规定，就像之前一样：一个子证明只有在它不位于某个在该行已经闭合了的其他子证明内部时，才能在该行被引用。那个灾难性的证明尝试违反了这条规定。第 3–4 行的子证明位于一个在第 5 行结束的子证明内部。因此它不能在第 7 行被引用。

这是我们必须排除的另一种情况：


\begin{fitchproof}
\hypo{a}{A}\PR
\open
	\hypo{b}{B}\AS
	\open
	\hypo{c}{C}\AS
	\have{bc}{B \eand C}\ai{b,c}
	\have{c2}{C}\ae{bc}
	\close
\close
\ifHTMLtarget
\have{bcab}{B \eif C}\by{不当的尝试以应用 $\eif$I}{b-c2}
\else
\have{bcab}{B \eif C}\by{不当的尝试}{}
	\have [\ ]{x}{} \by{以应用 $\eif$I}{b-c2}
\fi
\end{fitchproof}


这里我们试图引用一个从第 2 行开始到第 5 行结束的子证明——但第 5 行的语句不仅依赖于第 2 行的假设，还依赖于另一个我们在子证明结束时尚未解除的假设（第 3 行）。第 3 行开始的子证明在第 5 行时仍然是开放的。但是 $\eif$I 要求，被引用的子证明（即从第 2 行开始的子证明）的最后一行\emph{只能}依赖于该子证明本身的假设（以及它之前的任何东西），而不能依赖于它内部任何子证明的假设。具体来说，被引用的子证明的最后一行本身不能位于一个嵌套的子证明内部。

\factoidbox{在应用规则时引用一个子证明：

\begin{compactlist}
\item 被引用的子证明必须完整地位于应用该规则的那一步之前，
\item 被引用的子证明不能位于某个其他已关闭的子证明内部，且该子证明在引用行处已经关闭，并且
\item 被引用的子证明的最后一行不能出现在一个嵌套的子证明内部。
\end{compactlist}

}

最后一点用以强调规则的应用方式：如果一条规则要求引用单一行，就不能用子证明来替代；如果它要求引用子证明，也不能用单一行来替代。例如，下面这个证明是错误的：

\begin{fitchproof}
\hypo{a}{A}\PR
\open
	\hypo{b}{B}\AS
	\open
	\hypo{c}{C}\AS
	\have{bc}{B \eand C}\ai{b,c}
	\have{c2}{C}\ae{bc}
	\close
\ifHTMLtarget
	\have{c3}{C}\by{不当尝试以调用 R}{c-c2}
\else
\have{c3}{C}\by{不当尝试}{}
\have [\ ]{x}{} \by{以调用 R}{c-c2}
\fi
\close
\have[7]{bcab}{B \eif C}\ci{b-c3}
\end{fitchproof}

这里，我们试图在第6行通过重述规则R来证明$C$，但却引用了第3–5行的子证明。该子证明已经关闭，原则上可以在第6行引用它。（例如，我们可以用它通过$\eif$I来证明$C \eif C$。）但是重述规则R要求引用单一行，因此引用整个子证明是不允许的（即使该子证明包含我们想要重述的语句$C$）。在第6行引用包含$C$的单行（第3行或第5行）当然也不允许，因为这些行位于在第6行之前已经关闭的子证明中。

总是可以用任何假设开启一个子证明。然而，选择有用的假设需要一些策略。以一个任意的、古怪的假设开启子证明只会浪费证明的行数。例如，为了通过{\eif}I得到一个条件句，你必须在子证明中假设该条件句的前件。

同样，任何时候都可以关闭一个子证明（并解除其假设）。然而，在你得到有用结论之前，这样做并无帮助。一旦子证明关闭，在后续证明中只能引用整个子证明。反过来，那些要求引用子证明的规则，也要求子证明的最后一行是某种特定形式的语句。例如，只有当你所要证明的行的形式是$\metav{A} \eif \metav{B}$，且$\metav{A}$是你的子证明的假设，$\metav{B}$是你的子证明的最后一行时，才允许引用该子证明来应用~$\eif$I规则。

\section{双条件}

双条件的规则类似于条件规则的双重版本。

例如，为了证明“$F \eiff G$”，你必须能够在假设“$F$”下证明“$G$”\emph{并且}在假设“$G$”下证明“$F$”。因此，双条件引入规则（{\eiff}I）需要两个子证明。该规则的示意图如下： \factoidbox{


\begin{fitchproof}
	\open
		\hypo[i]{a1}{\metav{A}}\AS
		\have[j]{b1}{\metav{B}}
	\close
	\open
		\hypo[k]{b2}{\metav{B}}\AS
		\have[l]{a2}{\metav{A}}
	\close
	\have[\ ]{ab}{\metav{A}\eiff\metav{B}}\bi{a1-b1,b2-a2}
\end{fitchproof}

}
在$i$和$j$之间可以有任意多行，在$k$和$l$之间也可以有任意多行。而且，子证明可以按任意顺序出现，第二个子证明不必紧接在第一个之后。

双条件消去规则（{\eiff}E）允许你做比条件规则更多一点的操作。如果你有双条件句的左侧子句，你可以得到右侧子句。如果你有右侧子句，你可以得到左侧子句。因此我们允许：\factoidbox{


\begin{fitchproof}
	\have[m]{ab}{\metav{A}\eiff\metav{B}}
	\have[n]{a}{\metav{A}}
	\have[\ ]{b}{\metav{B}} \be{ab,a}
\end{fitchproof}

}
同样：\factoidbox{

\begin{fitchproof}
	\have[m]{ab}{\metav{A}\eiff\metav{B}}
	\have[n]{a}{\metav{B}}
	\have[\ ]{b}{\metav{A}} \be{ab,a}
\end{fitchproof}

}
注意，双条件式与其右半部分或左半部分可以分开引用，并且它们可以按任意顺序出现。然而，在引用~$\eiff$E 规则时，通常最好先引用双条件式本身。

\section{析取}
假设 Ludwig 是反动派。那么 Ludwig 要么是反动派，要么是自由意志主义者。毕竟，说 Ludwig 要么是反动派要么是自由意志主义者，比直接说 Ludwig 是反动派断定的内容更弱。

让我们强调一下这一点。假设 Ludwig 是反动派。由此可以推出，Ludwig \emph{要么}是反动派，\emph{要么}是一个金橘。同样，可以推出，\emph{要么} Ludwig 是反动派，\emph{要么}金橘是唯一的水果。同样，可以推出，\emph{要么} Ludwig 是反动派，\emph{要么}上帝死了。得出的许多结论听起来都很奇怪，但它们在\emph{逻辑上}没有问题（即使它们可能违反了各种隐含的会话准则）。

有了这些认识，我们给出析取引入规则：
\factoidbox{

\begin{fitchproof}
	\have[m]{a}{\metav{A}}
	\have[\ ]{ab}{\metav{A}\eor\metav{B}}\oi{a}
\end{fitchproof}

}
以及
\factoidbox{

\begin{fitchproof}
	\have[m]{a}{\metav{A}}
	\have[\ ]{ba}{\metav{B}\eor\metav{A}}\oi{a}
\end{fitchproof}

}
注意，这里的 \metav{B} 可以是\emph{任何}语句，因此下面这个证明完全可接受：

\begin{fitchproof}
	\hypo{m}{M}\PR
	\have{mmm}{M \eor ([(A\eiff B) \eif (C \eand D)] \eiff [E \eand F])}\oi{m}
\end{fitchproof}

如果要用真值表来验证这一点，将需要 128 行。

析取消去规则则稍微复杂一点。假设 Ludwig 要么是反动派，要么是自由意志主义者。你能得出什么结论？不能得出 Ludwig 是反动派；因为他也可能是自由意志主义者。同样，也不能得出 Ludwig 是自由意志主义者；因为他可能仅仅是反动派。单凭析取式本身，是难以直接用于有效推理的。

但是，假设我们能够以某种方式同时证明以下两点：第一，Ludwig 是反动派蕴涵着他是一位奥地利经济学家；第二，Ludwig 是自由意志主义者蕴涵着他是一位奥地利经济学家。那么，如果我们知道 Ludwig 要么是反动派要么是自由意志主义者，那么我们就知道，无论他是哪一种，Ludwig 都是一位奥地利经济学家。这一洞见可以表述为以下规则，即我们的析取消去（$\eor$E）规则：
\factoidbox{


\begin{fitchproof}
		\have[m]{ab}{\metav{A}\eor\metav{B}}
		\open
			\hypo[i]{a}{\metav{A}} \AS
			\have[j]{c1}{\metav{C}}
		\close
		\open
			\hypo[k]{b}{\metav{B}} \AS
			\have[l]{c2}{\metav{C}}
		\close
		\have[ ]{c}{\metav{C}}\oe{ab, a-c1,b-c2}
	\end{fitchproof}

}
显然，这比我们之前的规则写起来要稍微繁复一点，但原理相当简单。假设我们有一个析取式 $\metav{A} \eor \metav{B}$。假设我们有两个子证明，分别表明从假设~$\metav{A}$ 可以推出 $\metav{C}$，以及从假设~$\metav{B}$ 可以推出 $\metav{C}$。那么我们就可以推出 $\metav{C}$。和往常一样，在 $i$ 和 $j$ 之间、$k$ 和 $l$ 之间可以有很多行。此外，子证明和析取式可以以任何顺序出现，且不必相邻。

一些例子可能有助于说明这一点。考虑这个论证：
$$(P \eand Q) \eor (P \eand R) \therefore P$$
一个证明示例如下：


\begin{fitchproof}
		\hypo{prem}{(P \eand Q) \eor (P \eand R) }\PR
			\open
				\hypo{pq}{P \eand Q}\AS
				\have{p1}{P}\ae{pq}
			\close
			\open
				\hypo{pr}{P \eand R}\AS
				\have{p2}{P}\ae{pr}
			\close
		\have{con}{P}\oe{prem, pq-p1, pr-p2}
	\end{fitchproof}


下面是一个稍复杂一点的例子。考虑：
	$$ A \eand (B \eor C) \therefore (A \eand B) \eor (A \eand C)$$
以下是这个论证对应的一个证明：


\begin{fitchproof}
		\hypo{aboc}{A \eand (B \eor C)}\PR
		\have{a}{A}\ae{aboc}
		\have{boc}{B \eor C}\ae{aboc}
		\open
			\hypo{b}{B}\AS
			\have{ab}{A \eand B}\ai{a,b}
			\have{abo}{(A \eand B) \eor (A \eand C)}\oi{ab}
		\close
		\open
			\hypo{c}{C}\AS
			\have{ac}{A \eand C}\ai{a,c}
			\have{aco}{(A \eand B) \eor (A \eand C)}\oi{ac}
		\close
	\have{con}{(A \eand B) \eor (A \eand C)}\oe{boc, b-abo, c-aco}
	\end{fitchproof}


如果你一时难以自行构想出这个证明，也无需气馁。构造新颖证明的能力来源于练习，我们将在 \cref{s:stratTFL} 中介绍一些寻找证明的策略。现阶段的关键问题是，面对这个证明，你能否看出它符合我们制定的规则。这只需逐行检查，并确保每一行的证明依据都符合我们已确立的规则。

\section{矛盾与否定}\label{sec:contradiction}

我们只剩下一个联结词需要处理：否定。但要处理它，我们必须将否定与\emph{矛盾}联系起来。

一种有效的论证方法是使你的对手陷入自相矛盾。那时，你就使他们处于困境。他们必须放弃至少一个假设。我们将在证明系统中利用这一想法，通过添加一个新符号 “$\ered$” 到我们的证明中。这个符号应该读作类似“矛盾！”或“归谬！”或“但这太荒谬了！”。引入这个符号的规则是，每当我们明确地自相矛盾时，即每当我们在证明中发现一个语句及其否定同时出现时，就可以使用它：
\factoidbox{
\begin{fitchproof}
  \have[m]{na}{\enot\metav{A}}
  \have[n]{a}{\metav{A}}
  \have[ ]{bot}{\ered}\ne{na, a}
\end{fitchproof}
}
语句及其否定以何种顺序出现并不重要，它们也不必出现在相邻行上。按照惯例，先引用否定所在的行号，然后是它所否定的语句所在的行号。

显然，矛盾与否定之间联系紧密。规则 $\enot$E 让我们能从两个相互矛盾的语句——$\metav{A}$ 及其否定 $\enot \metav{A}$——进行到一个显式的矛盾 $\ered$。我们选择这个标签是有原因的：它是最基本的规则，让我们能从包含否定的前提（即 $\enot\metav{A}$）进行到一个不包含它的语句（即 $\ered$）。所以这是一条\emph{消除} $\enot$ 的规则。

我们说过 “$\ered$” 应该读作类似“矛盾！”的东西，但这并没有告诉我们太多关于这个符号本身的信息。大致有三种处理这个符号的方式：
\begin{compactlist}
    \item 我们可以把 “$\ered$” 看作真值函项逻辑的一个新的原子语句，但它只能取假值。
    \item 我们可以把 “$\ered$” 看作某个特定矛盾（例如 “$A \eand \enot A$”）的缩写。这将产生与上一种方式相同的效果——显然，“$A \eand \enot A$” 也只能取假值——但这意味着正式来说，我们不需要向真值函项逻辑添加新符号。
    \item 我们可以不把 “$\ered$” 看作真值函项逻辑的符号，而更像是我们证明中出现的一种\emph{标点符号}。（可以说，它与行号和竖线属于同一类。）
\end{compactlist}

第三种选择在哲学上非常吸引人，但在这里我们将\emph{正式}采用第一种。“$\ered$” 将被视为一个总为假的语句字母。这意味着我们可以在证明中像操作任何其他语句一样操作它。

我们仍然需要陈述否定引入的规则。规则非常简单：如果假设某事会导致矛盾，那么该假设必定是错误的。这一想法激发了以下规则：
\factoidbox{
\begin{fitchproof}
\open
    \hypo[i]{a}{\metav{A}}\AS
    \have[j]{nb}{\ered}
\close
\have[\ ]{na}{\enot\metav{A}}\ni{a-nb}
\end{fitchproof}
}
在 $i$ 行和 $j$ 行之间可以有你想要的任意多行。为了在实践中看到这一点，并观察其与否定的互动，考虑以下证明：
\begin{fitchproof}
    \hypo{d}{D}\AS
    \open
        \hypo{nd}{\enot D}\AS
        \have{ndr}{\ered}\ne{nd, d}
    \close
    \have{con}{\enot\enot D}\ni{nd-ndr}
\end{fitchproof}

如果假设 $\metav{A}$ 为真会导致矛盾，那么 $\metav{A}$ 不可能为真，即它必定为假，也就是说 $\enot\metav{A}$ 必定为真。当然，如果假设 $\metav{A}$ 为假（即假设 $\enot\metav{A}$ 为真）会导致矛盾，那么 $\metav{A}$ 不可能为假，即 $\metav{A}$ 必定为真。因此我们可以考虑以下规则：
\factoidbox{
\begin{fitchproof}
\open
    \hypo[i]{a}{\enot\metav{A}}\AS
    \have[j]{nb}{\ered}
\close
\have[\ ]{na}{\metav{A}}\ip{a-nb}
\end{fitchproof}
}
这条规则称为 \emph{间接证明}，因为它允许我们通过假设 $\metav{A}$ 的否定来间接地证明 $\metav{A}$。从形式上看，这条规则与 $\enot$I 非常相似，只是 $\metav{A}$ 和 $\enot\metav{A}$ 在假设与结论的位置上互换了。由于 $\enot\metav{A}$ 并非该规则的结论，我们并不是在引入 $\enot$，因此间接证明并不是一条引入任何联结词的规则。它也不是一条消去联结词的规则，因为它没有以 $\enot\metav{A}$ 为形式的独立前提，只有一个假设了 $\enot\metav{A}$ 这一形式的子证明。相比之下，$\enot$E 确实有一个形式为 $\enot\metav{A}$ 的前提：这就是为什么 $\enot$E 能消去 $\enot$，而间接证明却不能。\footnote{有些逻辑学家对间接证明心存疑虑，但对 $\enot$E 则没有异议。他们被称为“直觉主义者”。直觉主义者并不认同我们“每个语句要么为真要么为假”这一基本假设。他们认为 $\enot$ 的运作方式有所不同——对他们而言，从 $\metav{A}$ 推出 $\ered$ 能保证得到 $\enot \metav{A}$，但从 $\enot\metav{A}$ 推出 $\ered$ 并不能保证得到 $\metav{A}$，而只能得到 $\enot\enot\metav{A}$。因此，对他们来说，$\metav{A}$ 和 $\enot\enot\metav{A}$ 并不等价。}

利用 $\enot$I，我们能够从 $\metav{D}$ 给出 $\enot\enot\metav{D}$ 的证明。利用间接证明，我们可以反过来（用本质上相同的证明）进行推导。
\begin{fitchproof}
    \hypo{d}{\enot\enot D}\PR
    \open
        \hypo{nd}{\enot D}\AS
        \have{ndr}{\ered}\ne{d, nd}
    \close
    \have{con}{D}\ip{nd-ndr}
\end{fitchproof}

我们还需要最后一条规则。它是“$\ered$”的一种消去规则，称为 \emph{爆炸原理}。\footnote{这一原理的拉丁文名称是 \emph{ex contradictione quodlibet}，意为“从矛盾可推出任何命题”。} 如果我们得到了一个用“$\ered$”表示的矛盾，那么我们就可以推出任何我们想要的语句。作为一种论证规则，这如何能得到辩护呢？嗯，考虑一下英语中的一种修辞手法“……要是\emph{这}为真，我就吃了这顶帽子”。既然矛盾根本不可能是真的，那么如果有一个矛盾\emph{真的}为真，我不仅会吃掉我的帽子，我还会把它也吃掉。以下是该规则的形式表述：
\factoidbox{
\begin{fitchproof}
\have[m]{bot}{\ered}
\have[ ]{}{\metav{A}}\re{bot}
\end{fitchproof}
}
注意，\metav{A} 可以是\emph{任意}语句。

爆炸规则有点特别。看起来，\metav{A} 像魔术师从帽中取出兔子般凭空出现在我们的证明中。在寻找证明时，由于它看起来非常强大，我们很容易忍不住到处应用它。请克制住这种冲动：你\emph{只有已经得到}~$\ered$ 时才能应用它！而只有当你的假设导致矛盾时，你才能得到 $\ered$。

不过，从矛盾竟然能推出任何东西，这难道不奇怪吗？根据我们的蕴涵和有效性概念，并不奇怪。因为 \metav{A} 蕴涵 \metav{B} \ifeff{} 不存在任何对语句字母的赋值，能同时让 \metav{A} 为真且 \metav{B} 为假。而 $\ered$ 是一个矛盾——无论语句字母如何赋值，它都永远不会为真。既然没有任何赋值能使 $\ered$ 为真，当然也就没有任何赋值能使 $\ered$ 为真且 \metav{B} 为假！因此，根据我们的蕴涵定义，$\ered \entails \metav{B}$，无论 \metav{B} 是什么。矛盾蕴涵一切。\footnote{有些逻辑学家不接受这种观点。他们认为，如果 \metav{A} 蕴涵 \metav{B}，那么 \metav{A} 和 \metav{B} 之间必须有某种\emph{相关联系}——而 $\ered$ 和一个任意的句子~\metav{B} 之间并没有。因此，这些逻辑学家发展出其他“相关”逻辑，在那些逻辑里不允许使用爆炸规则。}

\emph{以上就是真值函项逻辑证明系统的所有基本规则。}

\practiceproblems

\problempart
下面两个“证明”是\emph{错误的}。请解释它们各自犯了什么错误。

\begin{fitchproof}
\hypo{abc}{(\enot L \eand A) \eor L}\PR
\open
\hypo{nla}{\enot L \eand A}\AS
\have{nl}{\enot L}\ae{nl}
    \have{a}{A}\ae{abc}
\close
\open
    \hypo{l}{L}\AS
    \have{red}{\ered}\ne{nl, l}
    \have{a2}{A}\re{red}
\close
\have{con}{A}\oe{abc, nla-a, l-a2}
\end{fitchproof}

\begin{fitchproof}
\hypo{abc}{A \eand (B \eand C)}\PR
\hypo{bcd}{(B \eor C) \eif D}\PR
\have{b}{B}\ae{abc}
\have{bc}{B \eor C}\oi{b}
\have{d}{D}\ce{bc, bcd}
\end{fitchproof}

\problempart
下面三个证明缺少引用（规则和行号）。请补充引用，将它们变成\emph{名副其实}的证明。此外，请写出每个证明对应的论证。

\begin{multicols}{2}
\begin{fitchproof}
\hypo{ps}{P \eand S}
\hypo{nsor}{S \eif R}
\have{p}{P}%\ae{ps}
\have{s}{S}%\ae{ps}
\have{r}{R}%\ce{nsor, s}
\have{re}{R \eor E}%\oi{r}
\end{fitchproof}

\begin{fitchproof}
\hypo{ad}{A \eif D}
\open
	\hypo{ab}{A \eand B}
	\have{a}{A}%\ae{ab}
	\have{d}{D}%\ce{ad, a}
	\have{de}{D \eor E}%\oi{d}
\close
\have{conc}{(A \eand B) \eif (D \eor E)}%\ci{ab-de}
\end{fitchproof}

\begin{fitchproof}
\hypo{nlcjol}{\enot L \eif (J \eor L)}
\hypo{nl}{\enot L}
\have{jol}{J \eor L}%\ce{nlcjol, nl}
\open
	\hypo{j}{J}
	\have{jj}{J \eand J}%\ai{j}
	\have{j2}{J}%\ae{jj}
\close
\open
	\hypo{l}{L}
	\have{red}{\ered}%\ne{nl, l}
	\have{j3}{J}%\re{red}
\close
\have{conc}{J}%\oe{jol, j-j2, l-j3}
\end{fitchproof}
\end{multicols}

\solutions
\problempart
\label{pr.solvedTFLproofs}
为以下每个论证各给出一个证明：


\begin{compactlist}
\item $J\eif\enot J \therefore \enot J$
\item $Q\eif(Q\eand\enot Q) \therefore \enot Q$
\item $A\eif (B\eif C) \therefore (A\eand B)\eif C$
\item $K\eand L \therefore K\eiff L$
\item $(C\eand D)\eor E \therefore E\eor D$
\item $A\eiff B, B\eiff C \therefore A\eiff C$
\item $\enot F\eif G, F\eif H \therefore G\eor H$
\item $(Z\eand K) \eor (K\eand M), K \eif D \therefore D$
\item $P \eand (Q\eor R), P\eif \enot R \therefore Q\eor E$
\item $S\eiff T \therefore S\eiff (T\eor S)$
\item $\enot (P \eif Q) \therefore \enot Q$
\item $\enot (P \eif Q) \therefore P$
\end{compactlist}

\chapter{构造证明}\label{s:stratTFL}

寻找证明并没有简单的诀窍，也没有什么可以替代练习。不过，这里有一些经验法则和策略可供参考。

\section{从目标倒推}

假设你正在尝试寻找某个结论~$\metav{C}$ 的证明，它将是证明的最后一行。首先要做的是观察~$\metav{C}$ 并思考其主逻辑算子的引入规则是什么。这能让你大致了解证明最后一行\emph{之前}应该是什么。该引入规则要求最后一行之前出现一个或两个其他语句，或者一个或两个子证明。此外，从~$\metav{C}$ 本身可以看出这些语句是什么，或者子证明的假设和结论是什么。然后，你就可以在最后一行之上写下那些语句，或者勾勒出子证明的框架，并将其作为新的目标。

例如：如果你的结论是一个条件句 $\metav{A}\eif\metav{B}$，就应该计划使用 {\eif}I 规则。这需要开启一个子证明，在其中假设~$\metav{A}$。该子证明应以~$\metav{B}$ 结束。接着，继续思考如何在该子证明内得到 $\metav{B}$，以及如何利用假设~$\metav{A}$。

如果你的目标是一个合取式、条件句或否定句，就应该以这种方式开始从目标倒推。接下来我们将详细描述在这几种情况下需要做什么。

\subsection*{从合取式倒推}

如果我们想证明 $\metav{A} \eand \metav{B}$，倒推意味着我们应在证明底部写下 $\metav{A} \eand \metav{B}$，并尝试用 $\eand$I 来证明它。在顶部，我们会写出证明的前提（如果有的话）。然后，在底部写下待证的语句。如果它是合取式，我们就用 $\eand$I 来证明。


\begin{fitchproof}
	\hypo{1}{\metav{P}_1}\PR
	\ellipsesline
	\hypo[k]{k}{\metav{P}_k}\PR
\ellipsesline
    \have[n]{n}{\metav{A}}{}
    \ellipsesline
	\have[m]{m}{\metav{B}}
    \have{4}{\metav{A} \eand \metav{B}}\ai{n,m}
  \end{fitchproof}


对于 $\eand$I，我们需要先证明 $\metav{A}$，再证明
$\metav{B}$。在最后一行，我们必须引用证明 $\metav{A}$ 和  $\metav{B}$ 的那些行（将来会填上），并使用~$\eand$I。
由垂直的 $\cdots$ 标出的部分仍有待填充。
我们暂用 $m$, $n$ 来标记这些行号。当证明完成时，这些占位符可被实际数字取代。

\subsection*{从条件句倒推}

如果目标是证明条件句 $\metav{A} \eif \metav{B}$，我们需要使用 $\eif$I。这要求一个以 $\metav{A}$ 开始、以~$\metav{B}$ 结束的子证明。我们这样来设置证明：


\begin{fitchproof}
\open
\hypo[n]{2}{\metav{A}}
\ellipsesline
\have[m]{3}{\metav{B}}
\close
\have{4}{\metav{A} \eif \metav{B}}\ci{2-3}
\end{fitchproof}


同样，我们会在行号处留下占位符。我们将最后一次推理记录为 $\eif$I，并引用该子证明。

\subsection*{从否定句倒推}

如果想证明 $\enot \metav{A}$，我们需要使用 $\enot$I。


\begin{fitchproof}
\open
\hypo[n]{2}{\metav{A}}
\ellipsesline
\have[m]{3}{\ered}
\close
\have{4}{\enot \metav{A}}\ni{2-3}
\end{fitchproof}


对于 $\enot$I，必须开启一个以 $\metav{A}$ 为假设的子证明；该子证明的最后一行必须是 $\ered$。我们将引用这个子证明，并应用~$\enot$I 规则。

在倒推时，应尽可能持续进行。所以，如果你为了证明 $\metav{A} \eif \metav{B}$ 而倒推，并已建立了一个要证明 $\metav{B}$ 的子证明，现在来看 $\metav{B}$。如果，比如说，它是一个合取式，就从它开始倒推，并将两个合取支写在子证明内。依此类推。

\subsection*{从析取式倒推}

当然，如果你的目标是析取式 $\metav{A} \eor \metav{B}$，也可以从它开始倒推。
$\eor$I 规则要求你拥有一个析取支，才能推出 $\metav{A} \eor \metav{B}$。
因此，倒推时，先选择一个析取支，从它（通过 $\eor$I 规则）推出 $\metav{A} \eor \metav{B}$，然后继续尝试证明你所选的析取支：


\begin{fitchproof}
	\ellipsesline
	\have[n]{2}{\metav{A}}
	\have{3}{\metav{A} \eor \metav{B}}\oi{2}
\end{fitchproof}


然而，你可能无法证明所选的那个析取支。此时就必须回溯。当你无法填满由垂直 $\cdots$ 标示的部分时，删除所有内容，尝试另一个析取支：


\begin{fitchproof}
	\ellipsesline
	\have[n]{2}{\metav{B}}
	\have{3}{\metav{A} \eor \metav{B}}\oi{2}
\end{fitchproof}


显然，全部删除并重新开始令人沮丧，应尽量避免。因此，如果你的目标是析取式，就\emph{不应}从倒推开始：应先尝试向前工作，仅当向前工作（以及使用 $\eand$I、$\eif$I 和 $\enot$I 的倒推）都无法进行时，才运用 $\eor$I 策略。

\section{向前工作}

你的证明可能包含前提。而且，如果你曾为证明一个条件句或否定句而进行过倒推，那么你已经建立了带有假设的子证明，并寻求在子证明中证明某个最终语句。这些前提和假设是你进行前向推理以填补证明中缺失步骤的出发点。这意味着要应用这些语句主算子的消去规则。规则的形式会告诉你该做什么。

\subsection*{从合取式向前工作}

要从形如 $\metav{A} \eand \metav{B}$ 的语句向前工作，我们使用 $\eand$E。该规则允许我们做两件事：推出 $\metav{A}$，以及推出 $\metav{B}$。因此，在已有 $\metav{A} \eand \metav{B}$ 的证明中，可通过在合取式下方紧接着写出 $\metav{A}$ 和/或 $\metav{B}$ 来进行向前工作：


\begin{fitchproof}
  \have[n]{1}{\metav{A} \eand \metav{B}}
  \have{2}{\metav{A}}\ae{1}
  \have{3}{\metav{B}}\ae{1}
\end{fitchproof}


通常在你所处的具体情境中，需要 $\metav{A}$ 和 $\metav{B}$ 中的哪一个是很清楚的。不过，把两者都写下来也无妨。

\subsection*{从析取式向前工作}

从析取式向前工作方式略有不同。要使用析取式，需运用 $\eor$E 规则。要应用此规则，仅知道待使用的析取式 $\metav{A} \eor \metav{B}$ 的析取支是什么是不够的，还必须牢记要证明什么。假设要证明~$\metav{C}$，且有 $\metav{A} \eor B$ 可用（该析取式可能是证明的前提、子证明的假设，或是已证的结论）。为应用 $\eor$E 规则，必须建立两个子证明：


\begin{fitchproof}
	\have[n]{1}{\metav{A} \eor \metav{B}}
	\open
	\hypo{2}{\metav{A}}
	\ellipsesline
	\have[m]{3}{\metav{C}}
	\close
	\open
	\hypo{4}{\metav{B}}
	\ellipsesline
	\have[k]{5}{\metav{C}}
	\close
	\have{6}{\metav{C}}\oe{1,(2)-3,(4)-5}
\end{fitchproof}


第一个子证明以第一个析取支 $\metav{A}$ 开始，以目标语句 $\metav{C}$ 结束。第二个子证明以另一个析取支 $\metav{B}$ 开始，同样以目标语句 $\metav{C}$ 结束。每个子证明都需进一步填充。然后，可以通过 $\eor$E 来证成目标语句 $\metav{C}$，并引用包含 $\metav{A} \eor \metav{B}$ 的行以及两个子证明。

\subsection*{从条件句向前工作}

要使用条件句 $\metav{A} \eif \metav{B}$，还需要前件 $\metav{A}$ 才能应用~$\eif$E。因此，从条件句向前工作时，将推出 $\metav{B}$，用 $\eif$E 证成它，并将 $\metav{A}$ 设为一个新的子目标：


\begin{fitchproof}
	\have[n]{1}{\metav{A} \eif \metav{B}}
	\ellipsesline
	\have[m]{2}{\metav{A}}
	\have{3}{\metav{B}}\ce{1,2}
\end{fitchproof}

\subsection*{从否定式向前工作}

最后，要使用一个否定式 $\enot \metav{A}$，你会应用 $\enot$E 规则。该规则除了需要 $\enot \metav{A}$，还需要对应的不带否定的语句~$\metav{A}$。由此得到的语句总是 $\ered$。因此，从否定式出发进行正向推理，在你打算用于 $\enot$I（或 IP）规则的子证明中尤为有效。当你已拥有 $\enot \metav{A}$ 并想要证明 $\ered$ 时，便可以从 $\enot \metav{A}$ 出发进行正向推理。做法是将 $\metav{A}$ 设为一个新的子目标。

\begin{fitchproof}
	\have[n]{1}{\enot \metav{A}}
	\ellipsesline
	\have[m]{2}{\metav{A}}
	\have{3}{\ered}\ne{1,2}
\end{fitchproof}

\section{策略实战}

假设我们想要证明论证 $(A \eand B) \eor (A \eand C) \therefore A \eand (B \eor C)$ 是有效的。我们通过写下前提和结论来开始证明。（在纸上，你应在两者之间留出尽可能多的空间，例如将前提写在页面上方，结论写在页面下方。）


\begin{fitchproof}
   \hypo{1}{(A \eand B) \eor (A \eand C)}\PR
\ellipsesline
  \have[n]{2}{A \eand (B \eor C)}
\end{fitchproof}


我们现在有两个选择：要么从结论出发进行反向推理，要么从前提出发进行正向推理。我们选取第二种策略：利用第 $1$ 行的析取式，并为 $\eor$E 规则建立所需的子证明。第 $1$ 行的析取式有两个析取支：$A \eand B$ 和 $A \eand C$。要证明的目标语句是 $A \eand (B \eor C)$。因此，这里需要建立两个子证明：一个以 $A \eand B$ 为假设，以 $A \eand (B \eor C)$ 为最后一行；另一个以 $A \eand C$ 为假设，也以 $A \eand (B \eor C)$ 为最后一行。第 $n$ 行结论的证明将基于 $\eor$E 规则，引用第 $1$ 行的析取式和这两个子证明。此时证明框架如下：


\begin{fitchproof}
	\hypo{1}{(A \eand B) \eor (A \eand C)}\PR
	\open
	\hypo{2}{A \eand B}\AS
	\ellipsesline
	\have[n]{6}{A \eand (B \eor C)}
	\close
	\open
	\hypo{7}{A \eand C}\AS
	\ellipsesline
	\have[m]{11}{A \eand (B \eor C)}
	\close
	\have{12}{A \eand (B \eor C)}\oe{1,2-6,(7)-11}
\end{fitchproof}


现在你面临两个独立的任务：分别完成这两个子证明。在第一个子证明中，我们从结论 $A \eand (B \eor C)$ 出发进行反向推理。这是一个合取式，因此在子证明内部，我们有两个独立的子目标：证明 $A$ 和证明 $B \eor C$。达成这些子目标后，就能通过 $\eand$I 规则证明第 $n$ 行。此时证明框架变为：


\begin{fitchproof}
	\hypo{1}{(A \eand B) \eor (A \eand C)}\PR
	\open
	\hypo{2}{A \eand B}\AS
	\ellipsesline
	\have[i]{4}{A}
	\ellipsesline
	\have[n][-1]{5}{B \eor C}
	\have[n]{6}{A \eand (B \eor C)}\ai{4,5}
	\close
	\open
	\hypo{7}{A \eand C}\AS
	\ellipsesline
	\have[m]{11}{A \eand (B \eor C)}
	\close
	\have{12}{A \eand (B \eor C)}\oe{1,2-6,(7)-11}
\end{fitchproof}


我们立刻可以看到，由第 $2$ 行通过 $\eand$E 规则即可得到第 $i$ 行。因此第 $i$ 行实际就是第 $3$ 行，可由第 $2$ 行根据 $\eand$E 规则证明。另一个子目标 $B \eor C$ 是一个析取式。我们将对第 $n-1$ 行运用“从析取式出发进行反向推理”的策略。我们需选择 $B$ 或 $C$ 其中之一作为子目标。选择 $C$ 行不通，最终将导致回溯。而若选择 $B$ 作为子目标，我们便可以通过对第 $2$ 行的合取式 $A \eand B$ 进行正向推理来得到它。因此，我们可以如下完成第一个子证明：


\begin{fitchproof}
	\hypo{1}{(A \eand B) \eor (A \eand C)}\PR
	\open
	\hypo{2}{A \eand B}\AS
	\have{3}{A}\ae{2}
	\have{4}{B}\ae{2}
	\have{5}{B \eor C}\oi{4}
	\have{6}{A \eand (B \eor C)}\ai{3,5}
	\close
	\open
	\hypo{7}{A \eand C}\AS
	\ellipsesline
	\have[m]{11}{A \eand (B \eor C)}
	\close
	\have{12}{A \eand (B \eor C)}\oe{1,2-6,7-11}
\end{fitchproof}


与第 $3$ 行类似，我们由第 $2$ 行通过 $\eand$E 规则得到第 $4$ 行。第 $5$ 行由第 $4$ 行根据 $\eor$I 规则证明，因为我们当时正在对一个析取式进行反向推理。

第一个子证明至此完成。第二个子证明几乎完全相同，我们将其留作练习。

请记住，开始时我们有两种选择：从前提出发进行正向推理，或从结论出发进行反向推理，我们选择了第一种。第二种选择也能导出一个证明，但形式会有所不同。最初的步骤应是从结论出发进行反向推理，设定两个子目标 $A$ 和 $B \eor C$，然后从前提出发进行正向推理来证明它们，例如：

\begin{fitchproof}
	\hypo{1}{(A \eand B) \eor (A \eand C)}\PR
	\open
	\hypo{2}{A \eand B}\AS
	\ellipsesline
	\have[k]{3}{A}
	\close
	\open
	\hypo{4}{A \eand C}\AS
	\ellipsesline
	\have[n][-1]{5}{A}
	\close
	\have{6}{A}\oe{1,2-3,(4)-(5)}
	\open
	\hypo{7}{A \eand B}\AS
	\ellipsesline
	\have[l]{8}{B \eor C}
	\close
	\open
	\hypo{9}{A \eand C}\AS
	\ellipsesline
	\have[m][-1]{10}{B \eor C}
	\close
	\have{11}{B \eor C}\oe{1,(7)-8,(9)-(10)}
	\have{12}{A \eand (B \eor C)}\ai{6,11}
\end{fitchproof}

我们将由你补全以 $\vdots$ 标示的缺失部分。

我们再举一例，以说明如何运用策略处理条件句和否定。语句 $(A \eif B) \eif
(\enot B \eif \enot A)$ 是一个重言式。我们来看看能否在不借助任何前提的情况下，运用策略为它找到一个证明。我们首先在纸页底部写下这个语句。由于正向推理无从着手（没有可作为起点的前提），我们采用反向推理，并为此建立一个子证明，利用 $\eif$I 规则来得到我们想要的语句 $(A
\eif B) \eif (\enot B \eif \enot A)$。其假设必须是我们想要证明的条件句的前件，即 $A \eif B$，而它的最后一行必须是后件，$\enot B \eif \enot A$。


\begin{fitchproof}
\open
\hypo{1}{A \eif B}\AS
\ellipsesline
\have[n]{7}{\enot B \eif \enot A}
\close
\have{8}{(A \eif B) \eif (\enot B \eif \enot A)}\ci{1-7}
\end{fitchproof}


新目标 $\enot B \eif \enot A$ 本身也是一个条件句，因此我们继续进行反向推理，再建立另一个子证明：


\begin{fitchproof}
	\open
	\hypo{1}{A \eif B}\AS
	\open
	\hypo{2}{\enot B}\AS
	\ellipsesline
	\have[n][-1]{6}{\enot A}
	\close
	\have{7}{\enot B \eif \enot A}\ci{2-(6)}
	\close
	\have{8}{(A \eif B) \eif (\enot B \eif \enot A)}\ci{1-7}
\end{fitchproof}

对于 $\enot A$，我们再次进行反向推理。为此，我们考虑 $\enot$I 规则。它要求建立一个以 $A$ 为假设、以 $\ered$ 为最后一行的子证明。因此，现在的证明变为：


\begin{fitchproof}
	\open
	\hypo{1}{A \eif B}\AS
	\open
	\hypo{2}{\enot B}\AS
	\open\hypo{3}{A}
	\ellipsesline
	\have[n][-2]{5}{\ered}
	\close
	\have{6}{\enot A}\ni{3-(5)}
	\close
	\have{7}{\enot B \eif \enot A}\ci{2-(6)}
	\close
	\have{8}{(A \eif B) \eif (\enot B \eif \enot A)}\ci{1-7}
\end{fitchproof}

现在，我们的目标是证明 $\ered$。此前讨论如何从否定式进行正向推理时曾指出：$\enot$E 规则允许你证明 $\ered$，这正是最内层子证明的目标。因此，我们寻找一个可用于正向推理的否定式：即第 $2$ 行的 $\enot B$。这意味着我们需要在该子证明内部推出 $B$，因为 $\enot$E 规则不仅需要 $\enot B$（我们已有），还需要 $B$。而 $B$ 本身，则可以通过对 $A \eif B$ 进行正向推理得到，因为 $\eif$E 规则允许我们证明该条件句的后件 $B$。$\eif$E 规则所需的条件句前件 $A$ 也已具备（位于第 $3$ 行）。于是我们完成证明如下：


\begin{fitchproof}
	\open
	\hypo{1}{A \eif B}\AS
	\open
	\hypo{2}{\enot B}\AS
	\open\hypo{3}{A}
	\have{4}{B}\ce{1,3}
	\have{5}{\ered}\ne{2,4}
	\close
	\have{6}{\enot A}\ni{3-5}
	\close
	\have{7}{\enot B \eif \enot A}\ci{2-6}
	\close
	\have{8}{(A \eif B) \eif (\enot B \eif \enot A)}\ci{1-7}
\end{fitchproof}

\section{从 $\ered$ 正向推理}\label{sec:backred}

运用这些策略时，你有时会遇到能够证明 $\ered$ 的情形。利用爆炸规则，这允许你证明\emph{任何}语句。因此，$\ered$ 在证明中如同万能牌。例如，假设你想为论证 $A \eor B, \enot A \therefore B$ 提供一个证明。你搭建起证明框架，在页面上方第 $1$ 行和第 $2$ 行写上前提 $A \eor B$ 和 $\enot A$，在页面底部写下结论 $B$。$B$ 没有主逻辑算子，因此无法对其进行反向推理。相反，你必须从 $A \eor B$ 进行正向推理：这需要两个子证明，如下所示：


\begin{fitchproof}
	\hypo{1}{A \eor B}\PR
	\hypo{7}{\enot A}\PR
	\open
	\hypo{2}{A} \AS
	\ellipsesline
	\have[m]{3}{B}
	\close
	\open
	\hypo{4}{B}\AS
	\ellipsesline
	\have[k]{5}{B}
	\close
	\have{6}{B}\oe{1,2-3,(4)-5}
\end{fitchproof}


注意，第 $2$ 行有 $\enot A$，而第一个子证明的假设是 $A$。利用 $\enot$E 规则可得到 $\ered$，进而根据 X 规则可由该 $\ered$ 得到第一个子证明的结论 $B$。回想一下，你可以使用重复规则 R 来复述一个已有的语句。因此，我们的证明如下：


\begin{fitchproof}
	\hypo{1}{A \eor B}\PR
	\hypo{7}{\enot A}\PR
	\open
	\hypo{2}{A} \AS
	\have{8}{\ered}\ne{7,2}
	\have{3}{B}\re{8}
	\close
	\open
	\hypo{4}{B}\AS
	\have{5}{B}\by{R}{4}
	\close
	\have{6}{B}\oe{1,2-3,4-5}
\end{fitchproof}

\section{进行间接证明}

在很多情况下，前推与后拉的策略最终都能成功。但在某些情况下，它们会失效。如果无法用这些方法直接证明 $\metav{A}$，就改用间接证明。为此，需要设立一个子证明，在其中假设 $\enot\metav{A}$，并尝试在该子证明中推出 $\ered$。

\begin{fitchproof}
\open
\hypo[n]{2}{\enot \metav{A}}\AS
\ellipsesline
\have[m]{3}{\ered}
\close
\have{4}{\metav{A}}\ip{2-3}
\end{fitchproof}

这里，我们必须以假设 $\enot \metav{A}$ 开启一个子证明；子证明的最后一行必须是 $\ered$。我们将引用该子证明，并应用间接证明规则。在子证明中，我们现在有了一个新的假设（即第 $n$ 行）可供利用。

假设我们使用了间接证明策略，或是处于其他需要推出 $\ered$ 的情境中。有什么好的着手点呢？最明显的候选者当然是使用一个否定句，因为（如上所见）$\enot$E 规则总是能推出 $\ered$。如果按照上述方式设定证明，试图用间接证明来证明 \metav{A}，那么你的子证明就会以 $\enot \metav{A}$ 作为假设——于是，为了在子证明中从前提出发证明 $\ered$，你接下来会以 \metav{A} 作为子证明内的一个目标。如果使用这种间接证明策略，你将会发现自己处于如下情境：

\begin{fitchproof}
\open
\hypo[n]{2}{\enot \metav{A}}\AS
\ellipsesline
\have[m][-1]{3}{\metav{A}}
\have{4}{\ered}\ne{2,3}
\close
\have{5}{\metav{A}}\ip{2-4}
\end{fitchproof}

这看起来有些奇怪：我们原本想证明 $\metav{A}$，但各种策略都失败了；于是我们退而使用间接证明。然而现在我们发现自己陷入了同样的困境：我们再次需要证明 $\metav{A}$。但请注意，我们现在处于一个子证明\emph{内部}，并且在这个子证明中，我们多了一个先前没有的假设（$\enot \metav{A}$）可以利用。让我们来看一个例子。

\section{排中律的间接证明}\label{s:proofLEM}

语句 $A \eor \enot A$ 是一个重言式，称为\emph{排中律}（LEM），因为它概括了这样的思想：要么 $\metav{A}$ 为真，要么 $\enot \metav{A}$ 为真，不存在两者皆非的中间情况。\footnote{你有时会遇到逻辑学家或哲学家谈论“tertium non datur”（拉丁语，意为“无第三者”）。这和排中律是同一个原则；意思是“没有第三条路”。对间接证明有疑虑的逻辑学家对排中律也有疑虑。} 它应该可以在没有任何前提的情况下得到证明。但后拉策略失效了：要使用 $\eor$I 得到 $A \eor \enot A$，我们必须证明 $A$ 或 $\enot A$——同样是在没有任何前提的情况下。这两者都不是重言式，所以我们无法证明其中任何一个。前推策略也不起作用，因为没有任何前提可以推进。因此，唯一的选项就是间接证明。

\begin{fitchproof}
	\open
	\hypo{2}{\enot (A \eor \enot A)}\AS
	\ellipsesline
	\have[m]{8}{\ered}
	\close
	\have{9}{A \eor \enot A}\ip{2-8}
\end{fitchproof}

现在我们确实有前提可以推进了：假设 $\enot(A \eor \enot A)$。为了利用它，我们通过 $\enot$E 来推出 $\ered$，这需要引用第 $1$ 行的假设，以及相应的非否定语句 $A \eor \enot A$（而这个语句尚待证明）。

\begin{fitchproof}
	\open
	\hypo{2}{\enot (A \eor \enot A)}\AS
	\ellipsesline
	\have[m][-1]{7}{A \eor \enot A}
	\have{8}{\ered}\ne{2,7}
	\close
	\have{9}{A \eor \enot A}\ip{2-8}
\end{fitchproof}

一开始，通过 $\eor$I 后拉来证明 $A \eor\enot A$ 没有成功。但现在我们处于一个不同的情境：我们需要在子证明内部证明 $A \eor\enot A$。一般来说，当遇到新的目标时，我们应该回过头来从基本策略开始。在这种情况下，我们应首先尝试从析取式 $A \eor \enot A$ 后拉，即我们必须选择一个析取支并尝试证明它。我们选择 $\enot A$。这样我们就能够在第 $m - 1$ 行使用 $\eor$I 来推出 $A \eor \enot A$。然后，为了使用 $\enot$I 从后拉策略证明 $\enot A$，我们开启另一个子证明。该子证明必须以 $A$ 为假设，并以 $\ered$ 为最后一行。

\begin{fitchproof}
	\open
	\hypo{2}{\enot (A \eor \enot A)}\AS
	\open
	\hypo{3}{A}\AS
	\ellipsesline
	\have[m][-3]{5}{\ered}
	\close
	\have{6}{\enot A}\ni{3-(5)}
	\have{7}{A \eor \enot A}\oi{6}
	\have{8}{\ered}\ne{2,7}
	\close
	\have{9}{A \eor \enot A}\ip{2-8}
\end{fitchproof}

在这个新的子证明内部，我们同样需要为 $\ered$ 提供理由。最好的办法是从一个否定式前推；第 $1$ 行的 $\enot(A \eor \enot A)$ 是我们唯一可用的否定式。然而，相应的非否定语句 $A \eor \enot A$，却可以通过 $\eor$I 直接从 $A$（我们在第 $2$ 行有 $A$）得到。我们完整的证明如下：

\begin{fitchproof}
	\open
	\hypo{2}{\enot (A \eor \enot A)}\AS
	\open
	\hypo{3}{A}\AS
	\have{4}{A \eor \enot A}\oi{3}
	\have{5}{\ered}\ne{2,4}
	\close
	\have{6}{\enot A}\ni{3-5}
	\have{7}{A \eor \enot A}\oi{6}
	\have{8}{\ered}\ne{2,7}
	\close
	\have{9}{A \eor \enot A}\ip{2-8}
\end{fitchproof}

\practiceproblems

\problempart
使用策略，为以下每个论证找出证明：


\begin{compactlist}
\item $A \eif B, A \eif C \therefore A \eif (B \eand C)$
\item $(A \eand B) \eif C \therefore A \eif (B \eif C)$
\item $A \eif (B \eif C) \therefore (A \eif B) \eif (A \eif C)$
\item $A \eor (B \eand C) \therefore (A \eor B) \eand (A \eor C)$
\item $(A \eand B) \eor (A \eand C) \therefore A \eand (B \eor C)$
\item $A \eor B, A \eif C, B \eif D \therefore C \eor D$
\item $\enot A \lor \enot B \therefore \enot(A \eand B)$
\item $A \eand \enot B \therefore \enot(A \eif B)$
\end{compactlist}

\problempart
制定从 $\metav{A} \eiff \metav{B}$ 出发进行前推与后拉的策略。

\problempart
使用策略，为以下每个语句找出证明：


\begin{compactlist}
\item $\enot A \eif (A \eif \ered)$
\item $\enot(A \eand \enot A)$
\item $[(A \eif C) \eand (B \eif C)] \eif [(A \lor B) \eif C]$
\item $\enot(A \eif B) \eif (A \eand \enot B)$
\item $(\enot A \eor B) \eif (A \eif B)$
\end{compactlist}


由于这些是在无前提情况下对语句的证明，你将从证明的\emph{底部}——即目标语句本身——开始，且没有前提。

\problempart
使用策略，为以下每个论证和语句找出证明：


\begin{compactlist}
\item $\enot\enot A \eif A$
\item $\enot A \eif \enot B \therefore B \eif A$
\item $A \eif B \therefore \enot A \eor B$
\item $\enot(A \eand B) \eif (\enot A \eor \enot B)$
\item $A \eif (B \eor C) \therefore (A \eif B) \eor (A \eif C)$
\item $(A \eif B) \lor (B \eif A)$
\item $((A \eif B) \eif A) \eif A$
\end{compactlist}


这些都将需要使用间接证明策略。尤其是最后三个会相当困难！

\chapter{真值函项逻辑的额外规则}\label{s:Further}
在\cref{s:BasicTFL}中，我们介绍了真值函项逻辑证明系统的基本规则。在本节中，我们将为系统添加一些额外规则。扩展后的证明系统使用起来会稍微容易一些。（不过，在\cref{s:Derived}中我们会看到，严格来说它们并非\emph{必需}。）

% \section{Reiteration}
% The first additional rule is \emph{reiteration} (R). This just allows us to repeat ourselves:
% \factoidbox{\begin{fitchproof}
% 	\have[m]{a}{\metav{A}}
% 	\have[\ ]{b}{\metav{A}} \by{R}{a}
% \end{fitchproof}}
% Such a rule is obviously legitimate; we could have used it in our proof in \cref{sec:backred}:
% \begin{fitchproof}
% 	\hypo{1}{A \eor B}
% 	\hypo{7}{\enot A}
% 	\open
% 	\hypo{2}{A}
% 	\have{8}{\ered}\ne{7,2}
% 	\have{3}{B}\re{8}
% 	\close
% 	\open
% 	\hypo{4}{B}
% 	\have{5}{B}\by{R}{4}
% 	\close
% 	\have{6}{B}\oe{1,2-3,4-5}
% \end{fitchproof}
% This is a fairly typical use of the R rule.

\section{析取三段论}
如下是一个非常自然的论证形式。


\begin{earg}
		\item Elizabeth 要么在马萨诸塞州，要么在华盛顿特区。
		\item 她不在华盛顿特区。
		\item[\texttherefore] 她在马萨诸塞州。
	\end{earg}


这被称为\emph{析取三段论}。我们如下将其加入证明系统：
\factoidbox{

\begin{fitchproof}
	\have[m]{ab}{\metav{A} \eor \metav{B}}
	\have[n]{nb}{\enot \metav{A}}
	\have[\ ]{con}{\metav{B}}\by{DS}{ab, nb}
\end{fitchproof}

}
以及
\factoidbox{

\begin{fitchproof}
	\have[m]{ab}{\metav{A} \eor \metav{B}}
	\have[n]{nb}{\enot \metav{B}}
	\have[\ ]{con}{\metav{A}}\by{DS}{ab, nb}
\end{fitchproof}

}
通常，析取式和一个析取支的否定可以以任意顺序出现，且不必相邻。不过，我们在引用规则时总是先列出析取式。

\section{拒取式}
另一个有用的推理模式体现在以下论证中：

\begin{earg}
		\item 如果 Hilary 赢得了选举，那么她就在白宫。
		\item 她不在白宫。
		\item[\texttherefore] 她没有赢得选举。
	\end{earg}

这一推理模式称为\emph{拒取式}。对应的规则是：
\factoidbox{

\begin{fitchproof}
	\have[m]{ab}{\metav{A}\eif\metav{B}}
	\have[n]{a}{\enot\metav{B}}
	\have[\ ]{b}{\enot\metav{A}}\mt{ab,a}
\end{fitchproof}

}
和往常一样，前提可以以任意顺序出现，但我们在引用规则时总是先列出条件句。

\section{双重否定消除}
另一个有用的规则是\emph{双重否定消除}。它正如其名：
\factoidbox{

\begin{fitchproof}
		\have[m]{dna}{\enot \enot \metav{A}}
		\have[ ]{a}{\metav{A}}\dne{dna}
	\end{fitchproof}

}
如此规定的理由是：在自然语言中，双重否定往往可以相互抵消。

话虽如此，但需要注意，语境和强调有时会阻止它们相互抵消。例如：“Jane is not \emph{not} happy”。有理由认为，我们不能从中推出“Jane is happy”，因为第一句话的意思应该被理解为等同于“Jane is not \emph{un}happy”。这与“Jane 处于一种彻底的漠然状态”是相容的。和往常一样，转向真值函项逻辑迫使我们牺牲了英语表达的某些细微之处。

\section{排中律}

假设我们能够证明：如果外面阳光明媚，那么 Bill 会带伞（以防晒伤）。再假设我们也能够证明：如果外面不是阳光明媚，那么 Bill 也会带伞（以防下雨）。那么，天气并没有第三种可能。所以，\emph{无论天气如何}，Bill 都会带伞。

这种思路促生了如下规则：
\factoidbox{

\begin{fitchproof}
		\open
			\hypo[i]{a}{\metav{A}}\AS
			\have[j]{c1}{\metav{B}}
		\close
		\open
			\hypo[k]{b}{\enot\metav{A}}\AS
			\have[l]{c2}{\metav{B}}
		\close
		\have[\ ]{ab}{\metav{B}}\tnd{a-c1,b-c2}
	\end{fitchproof}

}
这条规则有时被称为\emph{排中律规则}。你可以将其视为对 $\metav{A} \eor \enot \metav{A}$ 应用 $\eor$E 规则。
在 $i$ 和 $j$ 之间、以及 $k$ 和 $l$ 之间，可以有任意多行。
此外，这两个子证明可以以任意顺序出现，第二个子证明无需紧跟在第一个之后。

要看看这条规则的实际应用，考虑一下：
	$$P \therefore (P \eand D) \eor (P \eand \enot D)$$
下面是该论证对应的证明：


\begin{fitchproof}
		\hypo{a}{P}\PR
		\open
			\hypo{b}{D}\AS
			\have{ab}{P \eand D}\ai{a, b}
			\have{abo}{(P \eand D) \eor (P \eand \enot D)}\oi{ab}
		\close
		\open
			\hypo{nb}{\enot D}\AS
			\have{anb}{P \eand \enot D}\ai{a, nb}
			\have{anbo}{(P \eand D) \eor (P \eand \enot D)}\oi{anb}
		\close
		\have{con}{(P \eand D) \eor (P \eand \enot D)}\tnd{b-abo, nb-anbo}
	\end{fitchproof}


来看另一个例子：


\begin{fitchproof}
	\hypo{ana}{A \eif \enot A}\PR
	\open
		\hypo{a}{A}\AS
		\have{na}{\enot A}\ce{ana, a}
	\close
	\open
		\hypo{na1}{\enot A}\AS
		\have{na2}{\enot A}\by{R}{na1}
	\close
	\have{na3}{\enot A}\tnd{a-na, na1-na2}
\end{fitchproof}

\section{德摩根规则}
我们最后补充的一组规则被称为德摩根律（以 Augustus De~Morgan 命名）。这些规则的形式从真值表来看应该已经很熟悉了。

第一条德摩根规则是：
\factoidbox{

\begin{fitchproof}
	\have[m]{ab}{\enot (\metav{A} \eand \metav{B})}
	\have[\ ]{dm}{\enot \metav{A} \eor \enot \metav{B}}\dem{ab}
\end{fitchproof}

}
第二条德摩根规则是第一条的逆：
\factoidbox{

\begin{fitchproof}
	\have[m]{ab}{\enot \metav{A} \eor \enot \metav{B}}
	\have[\ ]{dm}{\enot (\metav{A} \eand \metav{B})}\dem{ab}
\end{fitchproof}

}
第三条德摩根规则是第一条的\emph{对偶}：
\factoidbox{

\begin{fitchproof}
	\have[m]{ab}{\enot (\metav{A} \eor \metav{B})}
	\have[\ ]{dm}{\enot \metav{A} \eand \enot \metav{B}}\dem{ab}
\end{fitchproof}

}
第四条则是第三条的逆：
\factoidbox{

\begin{fitchproof}
	\have[m]{ab}{\enot \metav{A} \eand \enot \metav{B}}
	\have[\ ]{dm}{\enot (\metav{A} \eor \metav{B})}\dem{ab}
\end{fitchproof}

}
\emph{以上便是我们用于真值函项逻辑的证明系统的所有补充规则。}

\practiceproblems
\solutions
\problempart
\label{pr.justifyTFLproof}
以下证明缺少其标注（规则和行号）。在需要的地方补充它们：


\begin{compactlist}
\item\begin{fitchproof}
\hypo{1}{W \eif \enot B}
\hypo{2}{A \eand W}
\hypo{2b}{B \eor (J \eand K)}
\have{3}{W}{}
\have{4}{\enot B} {}
\have{5}{J \eand K} {}
\have{6}{K}{}
\end{fitchproof}
\item\begin{fitchproof}
\hypo{1}{L \eiff \enot O}
\hypo{2}{L \eor \enot O}
\open
	\hypo{a1}{\enot L}
	\have{a2}{\enot O}{}
	\have{a3}{L}{}
	\have{a4}{\ered}{}
\close
\have{3b}{\enot\enot L}{}
\have{3}{L}{}
\end{fitchproof}
\item\begin{fitchproof}
\hypo{1}{Z \eif (C \eand \enot N)}
\hypo{2}{\enot Z \eif (N \eand \enot C)}
\open
	\hypo{a1}{\enot(N \eor  C)}
	\have{a2}{\enot N \eand \enot C} {}
	\have{a6}{\enot N}{}
	\have{b4}{\enot C}{}
		\open
		\hypo{b1}{Z}
		\have{b2}{C \eand \enot N}{}
		\have{b3}{C}{}
		\have{red}{\ered}{}
	\close
	\have{a3}{\enot Z}{}
	\have{a4}{N \eand \enot C}{}
	\have{a5}{N}{}
	\have{a7}{\ered}{}
\close
\have{3b}{\enot\enot(N \eor C)}{}
\have{3}{N \eor C}{}
\end{fitchproof}
\end{compactlist}

\problempart
为下列每个论证提供一个证明：

\begin{compactlist}
\item $E\eor F$, $F\eor G$, $\enot F \therefore E \eand G$
\item $M\eor(N\eif M) \therefore \enot M \eif \enot N$
\item $(M \eor N) \eand (O \eor P)$, $N \eif P$, $\enot P \therefore M\eand O$
\item $(X\eand Y)\eor(X\eand Z)$, $\enot(X\eand D)$, $D\eor M \therefore M$
\end{compactlist}

\chapter{证明论概念}\label{s:ProofTheoreticConcepts}

在本章，我们将引入一些新词汇。如下表达式：
$$\metav{A}_1, \metav{A}_2, \ldots, \metav{A}_n \proves \metav{C}$$
表示存在一个以 $\metav{C}$ 结尾的证明，其未被解除的假设是 $\metav{A}_1, \metav{A}_2, \ldots, \metav{A}_n$。若我们想说 \emph{不存在} 从 $\metav{A}_1$, $\metav{A}_2$, \dots,~$\metav{A}_n$ 出发且以 $\metav{C}$ 结尾的证明，则写作：$$\metav{A}_1, \metav{A}_2, \ldots, \metav{A}_n \nproves \metav{C}$$

符号“$\proves$”称为\emph{单推出符}。需特别注意，这不是我们在 \cref{s:SemanticConcepts} 中引入、用以符号化蕴涵的\emph{双推出符}（“$\entails$”）。单推出符“$\proves$”关乎证明的存在性；双推出符“$\entails$”则关乎赋值（或在一阶逻辑中使用时称为“解释”）的存在性。\emph{它们是两个截然不同的概念。}

借助“$\proves$”符号，我们可以引入更多术语。若存在一个没有未被解除假设的、对 $\metav{A}$ 的证明，我们写作：${} \proves \metav{A}$。此时，称 $\metav{A}$ 为\define{定理}。
	\factoidbox{\label{def:syntactic_tautology_in_sl}
		$\metav{A}$ 是\define{定理} \ifeff{} $\proves \metav{A}$
	}
\newglossaryentry{theorem}
{
name=定理,
description={一个可以在没有任何前提的情况下被证明的语句}
}

作为示例，假设我们想证明“$\enot (A \eand \enot A)$”是定理。这意味着我们需要一个\emph{没有}未被解除假设的、对“$\enot(A \eand \enot A)$”的证明。由于欲证语句的主逻辑算子是否定号，我们通常会以一个\emph{子证明}开始：在该子证明中，假设“$A \eand \enot A$”，并导出矛盾。因此，证明如下：

\begin{fitchproof}
		\open
			\hypo{con}{A \eand \enot A}\AS
			\have{a}{A}\ae{con}
			\have{na}{\enot A}\ae{con}
			\have{red}{\ered}\ne{na, a}
		\close
		\have{lnc}{\enot (A \eand \enot A)}\ni{con-red}
	\end{fitchproof}

由此，我们在没有（未被解除的）假设的情况下证明了“$\enot (A \eand \enot A)$”。这个特定的定理有时被称为\emph{不矛盾律}的一个实例。

要证明某语句是定理，只需找到一个合适的证明。而要证明某语句\emph{不是}定理，则通常困难得多。这要求你不仅证明某些尝试失败了，更要证明\emph{任何}证明都是不可能的。即便你尝试了千百种方法都未能证出该语句，也可能只是证明过长或过于复杂，你未能发现而已。

下面是另一个新术语：
	\factoidbox{
		两个语句 \metav{A} 和 \metav{B} 是\define{可互推的} \ifeff{} 二者均可从对方推出；即，同时满足 $\metav{A}\proves\metav{B}$ 与 $\metav{B}\proves\metav{A}$。
	}

\newglossaryentry{interderivable}
{
  name=相互可推导性,
  text = 相互可推导的,
description={两个语句共有的性质，当且仅当存在一个推导，可以从其中任一个语句推出另一个}
}

与证明定理的情形类似，证明两个语句可互推相对容易：只需提供一对推导。证明两个语句\emph{不可}互推则困难得多，这与证明一个语句不是定理同样困难。

第三个相关术语如下：
	\factoidbox{
	语句集 $\metav{A}_1,\metav{A}_2,\ldots, \metav{A}_n$ 是\define{不一致的} \ifeff{} 可以从它们（作为前提）推出矛盾~$\ered$，即 $\metav{A}_1,\metav{A}_2,\ldots, \metav{A}_n \proves
	\ered$。若它们并非\define{不一致的}，则称之为\define{一致的}。
	}

\newglossaryentry{inconsistent}
{    name={不一致性},
  description={一个语句集是不一致的 \ifeff{} 可以从它们推出矛盾~$\ered$},
    text={不一致的}
}

\newglossaryentry{consistent}
{    name={一致性},
  description={语句集是一致的 \ifeff{} \emph{无法}从它们推出矛盾~$\ered$},
    text={一致的}
}

要证明一组语句不一致很容易：只需以它们为前提，推出矛盾~$\ered$。而要证明一组语句\emph{不是}（即是一致的）则困难得多。这需要的不仅仅是一两个证明；它要求证明某类证明是\emph{不可能}存在的。

\
\\
下表总结了各类情形：是只需提供一两个证明，还是必须考察所有可能的证明。

\begin{center}
\begin{tabular}{l l l}
%\cline{2-3}
 & \textbf{是} & \textbf{否}\\
 \hline
%\cline{2-3}
定理？ & 只需一个证明 & 需要所有可能的证明\\
不一致？ &  只需一个证明  & 需要所有可能的证明\\
等价？ & 需要两个证明 & 需要所有可能的证明\\
一致？ & 需要所有可能的证明 & 只需一个证明\\
\end{tabular}
\end{center}

\practiceproblems
\problempart
证明以下各句是定理：
\begin{compactlist}
\item $O \eif O$
\item $N \eor \enot N$
\item $J \eiff [J\eor (L\eand\enot L)]$
\item $((A \eif B) \eif A) \eif A$
\end{compactlist}

\problempart
提供证明以表明以下各项：
\begin{compactlist}
\item $C\eif(E\eand G), \enot C \eif G \proves G$
\item $M \eand (\enot N \eif \enot M) \proves (N \eand M) \eor \enot M$
\item $(Z\eand K)\eiff(Y\eand M), D\eand(D\eif M) \proves Y\eif Z$
\item $(W \eor X) \eor (Y \eor Z), X\eif Y, \enot Z \proves W\eor Y$
\end{compactlist}

\problempart
证明以下每对句子都是可互推的：
\begin{compactlist}
\item $R \eiff E$, $E \eiff R$
\item $G$, $\enot\enot\enot\enot G$
\item $T\eif S$, $\enot S \eif \enot T$
\item $U \eif I$, $\enot(U \eand \enot I)$
\item $\enot (C \eif D), C \eand \enot D$
\item $\enot G \eiff H$, $\enot(G \eiff H)$
\end{compactlist}

\problempart
如果你知道 $\metav{A}\proves\metav{B}$，那么对于 $(\metav{A}\eand\metav{C})\proves\metav{B}$ 你能说什么？对于 $(\metav{A}\eor\metav{C})\proves\metav{B}$ 呢？请解释你的答案。

\

\problempart 在本章中，我们声称要证明两个句子不可互推，其难度与证明一个句子不是定理是一样的。我们为何这样声称？ (\emph{提示}：考虑一个句子是定理 \ifeff{} \metav{A} 和 \metav{B} 可互推。)

\chapter{派生规则}\label{s:Derived}
在本节中，我们将了解为何将证明系统的规则分两批引入。具体来说，我们想要说明 \cref{s:Further} 中介绍的附加规则并非严格必要，它们可以从 \cref{s:BasicTFL} 的基本规则中派生出来。

\section{重述规则的推导}
为了说明从其他规则“派生”一条规则意味着什么，首先考虑重述规则。它是我们系统中的一条基本规则，但并非不可或缺。假设在你的推导中，某一行上有一个句子：
\begin{fitchproof}
	\have[m]{a}{\metav{A}}
\end{fitchproof}
现在你想在第 $k$ 行重复它。你可以直接调用重述规则（R）。但同样地，你也可以用 \cref{s:BasicTFL} 的其他基本规则来实现：
\begin{fitchproof}
	\have[m]{a}{\metav{A}}
	\have[k]{aa}{\metav{A} \eand \metav{A}}\ai{a, a}
	\have{a2}{\metav{A}}\ae{aa}
\end{fitchproof}
需要明确的是：这不是一个证明，而是一个证明\emph{图式}。毕竟，它使用的是元变元“$\metav{A}$”而不是真值函项逻辑的一个语句。但要点很简单：不论我们用什么真值函项逻辑语句代入“$\metav{A}$”，也不论我们当时正在处理哪几行，我们都能产生一个真正的证明。因此，你可以将其视为一个生成证明的配方。

事实上，这个配方告诉我们：任何使用重述规则（R）能证明的东西，我们都可以（多费一行）仅用 \cref{s:BasicTFL} 的其他基本规则（不使用 R）来证明。这就是我们所说“重述规则（R）可从其他基本规则派生”的含义：任何能用 R 论证的内容，都可以只用其他基本规则来论证。

\section{析取三段论规则的推导}
假设在一个证明中，你有如下形式的内容：
\begin{fitchproof}
	\have[m]{ab}{\metav{A}\eor\metav{B}}
	\have[n]{na}{\enot \metav{A}}
\end{fitchproof}
现在你想在第 $k$ 行证明 $\metav{B}$。你可以使用析取三段论（DS）规则（于 \cref{s:Further} 引入）。但同样地，你也可以用 \cref{s:BasicTFL} 的\emph{基本}规则来完成：
\begin{fitchproof}
		\have[m]{ab}{\metav{A}\eor\metav{B}}
		\have[n]{na}{\enot \metav{A}}
		\open
			\hypo[k]{a}{\metav{A}}\AS
			\have{red}{\ered}\ne{na, a}
			\have{b1}{\metav{B}}\re{red}
		\close
		\open
			\hypo{b}{\metav{B}}\AS
			\have{b2}{\metav{B}}\by{R}{b}
		\close
	\have{con}{\metav{B}}\oe{ab, a-b1, b-b2}
\end{fitchproof}
因此，析取三段论（DS）规则同样可以从我们的基本规则中派生出来。将其加入系统并未使任何新的证明成为可能。任何时候你使用 DS 规则，总可以通过额外的几行，仅用我们的基本规则来证明同样的东西。它是一个\emph{派生}规则。

\section{否定后件规则的推导}
假设你的证明中有如下内容：
\begin{fitchproof}
	\have[m]{ab}{\metav{A}\eif\metav{B}}
	\have[n]{a}{\enot\metav{B}}
\end{fitchproof}
现在你想在第 $k$ 行证明 $\enot \metav{A}$。你可以使用否定后件（MT）规则（于 \cref{s:Further} 引入）。同样地，你也可以用 \cref{s:BasicTFL} 的\emph{基本}规则来完成：
\begin{fitchproof}
	\have[m]{ab}{\metav{A}\eif\metav{B}}
	\have[n]{nb}{\enot\metav{B}}
		\open
		\hypo[k]{a}{\metav{A}}\AS
		\have{b}{\metav{B}}\ce{ab, a}
		\have{nb1}{\ered}\ne{nb, b}
		\close
	\have{no}{\enot\metav{A}}\ni{a-nb1}
\end{fitchproof}
同样，MT 规则也能从 \cref{s:BasicTFL} 的\emph{基本}规则中推导出来。

\section{双重否定消去规则的推导}
考虑下面的推导演示：
\begin{fitchproof}
	\have[m]{m}{\enot \enot \metav{A}}
	\open
		\hypo[k]{a}{\enot\metav{A}}\AS
		\have{a1}{\ered}\ne{m, a}
	\close
	\have{con}{\metav{A}}\ip{a-a1}
\end{fitchproof}
同样，我们可以从 \cref{s:BasicTFL} 的\emph{基本}规则推导出双重否定消去（DNE）规则。

\section{排中律的推导}
假设你想使用排中律（LEM）规则来证明某命题，即你的证明中有如下结构：
\begin{fitchproof}
  \open
  \hypo[m]{a}{\metav{A}}\AS
  \have[n]{aaa}{\metav{B}}
  \close
  \open
  \hypo[k]{b}{\enot\metav{A}}\AS
  \have[l]{bbb}{\metav{B}}
  \close
\end{fitchproof}
现在你想在第 $l+1$ 行证明 $\metav{B}$。来自 \cref{s:Further} 的 LEM 规则允许你这样做。但是，你能用 \cref{s:BasicTFL} 的\emph{基本}规则做到这一点吗？

一种选择是先证明 $\metav{A} \eor \enot\metav{A}$，然后应用析取消去规则（$\eor$E），即分情形证明：
\begin{fitchproof}
  \open
  \hypo[m]{a}{\metav{A}}\AS
  \have[n]{aaa}{\metav{B}}
  \close
  \open
  \hypo[k]{b}{\enot\metav{A}}\AS
  \have[l]{bbb}{\metav{B}}
  \close
  \ellipsesline
  \have[i]{tnd}{\metav{A} \eor \enot \metav{A}}
  \have[i+1]{fin}{\metav{B}}\oe{tnd, a-aaa,b-bbb}
\end{fitchproof}
（我们在 \cref{s:proofLEM} 中给出了一个仅使用基本规则对 $\metav{A} \eor \enot\metav{A}$ 的证明。）

下面是另一种方法，比前面的方法略显复杂。你需要做的是将原来的两个子证明嵌套在另一个子证明内部。该子证明的假设是 $\enot \metav{B}$，最后一行是 $\ered$。因此，整个子证明正是那种你可以使用间接证明（IP）来推出 $\metav{B}$ 的类型。在证明内部，你需要多做一点工作才能得到 $\ered$：
\begin{fitchproof}
  \open
  \hypo[m]{nb}{\enot\metav{B}}\AS
  \open
  \hypo{a}{\metav{A}}\AS
  \ellipsesline
  \have[n]{aaa}{\metav{B}}
  \have{aaaa}{\ered}\ne{nb, aaa}
  \close
  \open
  \hypo{b}{\enot\metav{A}}\AS
  \ellipsesline
  \have[l]{bbb}{\metav{B}}
  \have{bbbb}{\ered}\ne{nb, bbb}
  \close
  \have{na}{\enot\metav{A}}\ni{(a)-(aaaa)}
  \have{nna}{\enot\enot\metav{A}}\ni{(b)-(bbbb)}
  \have{bot}{\ered}\ne{nna, na}
  \close
  \have{B}{\metav{B}}\ip{nb-(bot)}
\end{fitchproof}
注意，由于我们在顶部增加了一个假设，并且在子证明内部增加了额外的结论，行号发生了变化。你可能需要花点时间理解这个证明的运作方式。

\section{德摩根规则的推导}
以下是我们如何推导第一条德摩根规则的演示：
\begin{fitchproof}
		\have[m]{nab}{\enot (\metav{A} \eand \metav{B})}
		\open
			\hypo[k]{a}{\metav{A}}\AS
			\open
				\hypo{b}{\metav{B}}\AS
				\have{ab}{\metav{A} \eand \metav{B}}\ai{a,b}
				\have{nab1}{\ered}\ne{nab, ab}
			\close
			\have{nb}{\enot \metav{B}}\ni{(b)-(nab1)}
			\have{dis}{\enot\metav{A} \eor \enot \metav{B}}\oi{nb}
		\close
		\open
			\hypo{na1}{\enot \metav{A}}\AS
			\have{dis1}{\enot\metav{A} \eor \enot \metav{B}}\oi{na1}
		\close
		\have{con}{\enot \metav{A} \eor \enot \metav{B}}\tnd{a-(dis), (na1)-(dis1)}
	\end{fitchproof}
以下是我们如何推导第二条德摩根规则的演示：
\begin{fitchproof}
		\have[m]{nab}{\enot \metav{A} \eor \enot \metav{B}}
		\open
			\hypo[k]{ab}{\metav{A} \eand \metav{B}}\AS
			\have{a}{\metav{A}}\ae{ab}
			\have{b}{\metav{B}}\ae{ab}
			\open
				\hypo{na}{\enot \metav{A}}\AS
				\have{c1}{\ered}\ne{na, a}
			\close
			\open
				\hypo{nb}{\enot \metav{B}}\AS
				\have{c2}{\ered}\ne{nb, b}
			\close
			\have{con2}{\ered}\oe{nab, (na)-(c1), (nb)-(c2)}
		\close
		\have{nab}{\enot (\metav{A} \eand \metav{B})}\ni{ab-(con2)}
	\end{fitchproof}
类似的演示可以用来解释如何推导第三和第四条德摩根规则。这些留作练习。

\practiceproblems

\problempart
给出证明图式，以证明确立第三和第四条德摩根规则作为派生规则的合理性。

\

\problempart
你在回答 \crefrange{s:Further}{s:ProofTheoreticConcepts} 的练习题时提供的证明使用了派生规则。请在这些证明中，将派生规则的使用替换为仅使用基本规则。你会在得到的结果证明中发现一些“重复”；在这些情况下，请提供一个仅使用基本规则的简化证明。（这将让你体会到派生规则的力量，以及所有规则是如何相互作用的。）

\

\problempart
给出 $\metav{A} \eor \enot\metav{A}$ 的一个证明。然后给出一个\emph{仅使用基本规则}的证明。

\

\problempart
证明：如果你将排中律（LEM）作为一条基本规则，则可以确立间接证明（IP）规则作为一条派生规则的合理性。也就是说，假设你有如下证明：

\begin{fitchproof}
  \open
  \hypo[m]{a}{\enot\metav{A}}\AS
  \have[\ ]{aa}{\dots}
  \have[n]{aaa}{\ered}
  \close
\end{fitchproof}

如何在不使用 IP 但使用 LEM 及所有其他基本规则的情况下，利用以上推导结构来证明 \metav{A}？

\

\problempart
仅使用基本规则证明第一条德摩根律，且\emph{不得使用 LEM}。（当然，你可以将使用 LEM 的证明与 LEM 本身的证明结合起来。请尝试直接给出一个证明。）

\chapter{可靠性与完备性}
\label{sec:soundness_and_completeness}

在 \cref{s:ProofTheoreticConcepts} 中，我们看到可以利用推导来检验一些原本用真值表检验的性质。我们不仅能证明一个结论可以从某些前提中推导出来，还能用推导来判断一个语句是否为定理，或者两个语句是否可互推。我们也开始以相同的方式使用单竖线（$\proves$）和双竖线（$\entails$）：若能通过真值表证明 \metav{A} 是重言式，则写作 $\entails \metav{A}$；若能通过推导证明 \metav{A} 是定理，则写作 $\proves \metav{A}$。

此时你或许会想，这两种竖线是否总以相同的方式运作：若能通过真值表证明 \metav{A} 是重言式，是否也总能通过推导证明它是定理？反之是否亦然？对于有效论证以及逻辑等值的语句对，这些关系是否也成立？事实证明，所有这些及更多类似问题的答案都是肯定的。我们可以通过证明以下这一点来确立它：基于真值表的概念与基于推导的相应概念是等价的。例如，对于任意语句 $\metav{A}$，我们都能证明 $\entails \metav{A}$ \ifeff{} $\proves \metav{A}$。

首先，我们需要分别针对真值表和推导来定义所有的逻辑概念。其中大部分工作已经完成。所有真值表相关的定义均在 \cref{s:SemanticConcepts} 中给出。我们也已经给出了定理以及可互推语句对的证明论定义。其他定义随之自然可得。对于大多数逻辑性质，我们都可以设计使用推导的检验方法；而那些无法直接检验的性质，则可以通过我们已经能定义的概念来定义。

例如，我们将定理定义为无需任何前提即可推导出的语句（第~\pageref{def:syntactic_tautology_in_sl} 页）。我们可以将 \define{真值函项逻辑中的不一致语句}\label{def:syntactic_contradiction_in_sl} 定义为可以从中推导出矛盾“$\ered$”的语句。\footnote{注意，$\metav{A} \proves \ered$ \ifeff{} $\lnot \metav{A}$ 是定理。}偶然语句的句法定义则略有不同。我们并不像使用真值表那样，拥有一种实用的、有限的方法来通过推导证明一个语句是偶然的。因此，我们只能通过否定其他概念来定义“偶然语句”：一个语句是 \define{真值函项逻辑中的证明论偶然语句} \label{def:syntactically_contingent_in_sl}，当且仅当它既不是定理也不是不一致语句。

一组语句是 \define{真值函项逻辑中的不一致集合}\label{def:syntactically_inconsistent_ in_sl}，当且仅当可以从它们推导出矛盾 $\ered$。另一方面，与偶然性类似，我们也没有一种实用的有限方法来直接检验一致性。因此，我们再次需要通过否定方式定义一个术语：一组语句是 \define{真值函项逻辑中的一致集合}\label{def:syntactically consistent in SL}，当且仅当它不是不一致的。

最后，一个论证是 \define{真值函项逻辑中可证有效的} \label{def:syntactically_valid_in_SL}，当且仅当存在一个从其前提推出其结论的推导。所有这些定义均在 \cref{table:truth_tables_or_derivations} 中给出。

\def\tmptable{\textbf{概念} 		&	\textbf{真值表（语义）定义} 	&	\textbf{证明论（句法）定义} \\ \hline \hline
重言式/定理   &	主逻辑算子下的真值列全为 T 的语句 & 可以不依赖任何前提而推导出的语句。	 \\ \hline
矛盾式/不一致语句		&	主逻辑算子下的真值列全为 F 的语句  &	可以从中推导出矛盾~$\ered$ 的语句\\ \hline
偶然语句	&	主逻辑算子下的真值列既包含 T 也包含 F 的语句 & 既不是定理也不是不一致的语句 \\ \hline
逻辑等价的/可互推的语句 &	主逻辑算子下的真值列完全相同。& 语句之间可以相互推导出来	\\ \hline
不可满足的/不一致的语句集	&	在真值表中没有一行能使它们全部为真的语句集。	& 可以从中推导出矛盾~$\ered$ 的语句集 \\ \hline
可满足的/一致的语句集	&	在真值表中至少有一行能使它们全部为真的语句集。 & 无法从中推导出矛盾~$\ered$ 的语句集	\\ \hline
有效论证		&	在其真值表中，没有这样一行：所有前提的主逻辑算子下为 T 而结论的主逻辑算子下为 F。  & 可以从前提推导出结论的论证。}

\ifHTMLtarget


\begin{table}
\begin{tabularx}{\textwidth}{X||X|X}
\tmptable
\end{tabularx}
\caption{定义逻辑概念的两种方式。}
\label{table:truth_tables_or_derivations}
\end{table}


\else


\begin{sidewaystable}\small
\tabulinesep=1ex
\begin{tabu}{X[.5,c,m] ||X[1,l,m] |X[1,l,m]}
\tmptable
\end{tabu}
\caption{定义逻辑概念的两种方式。}
\label{table:truth_tables_or_derivations}
\end{sidewaystable}


\fi

我们现在有了成对的概念，它们分别从语义（使用赋值和真值表）和证明论（基于自然演绎）的角度各定义了一次。如何确保这两种定义总是一致的？完整的证明远超出本书的范围。不过，我们可以勾勒其大致轮廓。我们将着重证明两个有效性概念是等价的。由此，其他概念的等价性便可迅速得出。证明必须从两个方向进行。首先，我们必须证明在证明论上有效的论证在语义上也是有效的。换句话说，凡是能用推导证明的，也能用真值表证明。用符号表示，即我们要证明 $\proves$ 蕴涵 $\entails$。随后，我们需要证明另一个方向：$\entails$ 蕴涵 $\proves$。

\newglossaryentry{soundness}
{
name=可靠性,
description={一个逻辑系统具有的性质，当且仅当 $\proves $ 蕴涵 $\entails $}
}

从 $\proves $ 到 $\entails $ 的论证，就是 \define{\gls{soundness}} 问题。\label{def:soundness} 一个证明系统是 \define{可靠的（sound）}，如果不存在这样的论证推导：该论证的真值表可以显示它是无效的。\label{def_Soundness} 要证明该系统是可靠的，就需要证明\emph{任何}可能的证明都是一个有效论证的证明。仅仅在许多有效论证上尝试证明成功，而在无效论证上尝试证明失败是不够的。这个证明本身需要一些细致的考量。简而言之，它需要证明每一条推理规则都保持如下性质：其结论被该证明的前提连同所有开放的假设所蕴涵。你可以在 \cref{ch:Soundness} 找到这个证明的详细过程。

可靠性意味着 $\metav{A} \proves \metav{B}$ 蕴涵 $\metav{A}
\entails \metav{B}.$ 那么另一个方向呢？为什么我们应当认为\emph{所有}能用真值表证明为有效的论证，也都能用推导来证明？

\newglossaryentry{completeness}
{
name=完备性,
description={一个逻辑系统具有的性质，当且仅当 $\entails $ 蕴涵 $\proves $}
}

这就是完备性问题。一个证明系统具有 \define{\gls{completeness}}\label{def:completeness} 性质，当且仅当每个语义上有效的论证都存在一个推导。证明一个系统是完备的，通常比证明它可靠更难。证明一个系统可靠，本质上就是表明证明系统中的所有规则都按预期工作。而证明一个系统完备，则意味着要表明系统中包含了\emph{所有}必需的规则，没有任何遗漏。证明这一点超出了本书的范围。重要的是，幸运的是，真值函项逻辑的证明系统既是可靠的，又是完备的。并非所有的证明系统或形式语言都如此。正因为真值函项逻辑如此，我们可以根据当前任务，选择给出证明或真值表——哪种更容易就用哪种。

既然我们知道真值表方法与推导方法可以互换，你就可以为任意给定问题选择想使用的方法。学生通常更喜欢使用真值表，因为它们可以纯机械地生成，这看起来“更容易”。然而，我们已经看到，真值表仅在涉及几个语句字母后就会变得大得不可思议。另一方面，有些情况下根本无法使用证明。我们在语法上将偶然语句定义为不能被证明是重言式或矛盾式的语句。证明这类否定性陈述没有实际可行的方法。我们永远无法确知，是否存在某个证明能表明一个语句是矛盾式，只是我们尚未找到。这种情况下，我们别无选择，只能诉诸真值表。类似地，我们可以用推导来证明两个语句等价，但如果我们想证明它们\emph{不}等价呢？我们无法证明自己永远找不到相关的推导。于是我们又不得不求助于真值表。

\Cref{table.ProofOrModel} 总结了何时最好给出证明，何时最好给出真值表。

\def\tmptable{\textbf{逻辑性质} 	&	\textbf{证明其成立} 	&	\textbf{证明其不成立} \\ \hline \hline
为重言式或定理 		& 推导该语句  						& 在该语句的真值表中找到一个值为假的赋值行 \\ \hline
为矛盾式 	&  从该语句推导出 $\ered$  		 & 在该语句的真值表中找到一个值为真的赋值行\\ \hline
具偶然性 			& 在该语句的真值表中分别找到一个值为假和一个值为真的赋值行 & 推导该语句或其否定\\ \hline
具等价性 			& 从一个语句推导出另一个（双向推导） 		 & 在两个语句的真值表中找到一个它们真值不同的赋值行\\ \hline
具一致性 		& 在所有语句的真值表中找到一个它们全部为真的赋值行 & 从这些语句推导出 $\ered$\\ \hline
具有效性 				& 从前提推导出结论 & 在真值表中找到一个赋值行，使得所有前提为真而结论为假。 \\ }

\ifHTMLtarget


\begin{table}
\begin{tabularx}{\textwidth}{X||X|X}
\tmptable
\end{tabularx}
\caption{何时提供真值表，何时提供证明。}
\label{table.ProofOrModel}
\end{table}


\else


\begin{table}\small
\tabulinesep=1ex
\begin{tabu}{X[.7,c,m] ||X[1,l,m] |X[1,l,m]}
\tmptable
\end{tabu}
\caption{何时提供真值表，何时提供证明。}
\label{table.ProofOrModel}
\end{table}


\fi

\practiceproblems
\noindent\problempart 对下列各项，使用推导或真值表之一完成。


\begin{compactlist}%[label=(\arabic*)]
\item 证明 $A \eif [((B \eand C) \eor D) \eif A]$ 是一个定理。
\item 证明 $A \eif (A \eif B)$ 不是定理。
\item 证明语句 $A \eif \enot{A}$ 不是矛盾式。
\item 证明语句 $A \eiff \enot A$ 是矛盾式。
\item 证明语句 $ \enot (W \eif (J \eor J)) $ 是偶然的。
\item 证明语句 $ \enot(X \eor (Y \eor Z)) \eor (X \eor (Y \eor Z))$ 不是偶然的。
 \item 证明语句 $B \eif \enot S$ 等价于语句 $\enot \enot B \eif \enot S$。
\item 证明语句 $ \enot (X \eor O) $ 不等价于语句 $X \eand O$。
\item 证明语句 $\enot(A \eor B)$, $C$, $C \eif A$ 是不一致的。
\item 证明语句 $\enot(A \eor B)$, $\enot{B}$, $B \eif A$ 是一致的。
\item 证明 $\enot(A \eor (B \eor C)) \therefore \enot{C}$ 是有效的。
\item 证明 $\enot(A \eand (B \eor C)) \therefore \enot{C}$ 是无效的。
\end{compactlist}

\noindent\problempart 对下列各项，使用推导或真值表之一完成。

\begin{compactlist}%[label=(\arabic*)]
\item 证明 $A \eif (B \eif A)$ 是定理。
\item 证明 $\enot (((N \eiff Q) \eor Q) \eor N)$ 不是定理。
\item 证明 $ Z \eor (\enot Z \eiff Z) $ 是偶然的。
\item 证明 $ (L \eiff ((N \eif N) \eif L)) \eor H $ 不是偶然的。
\item 证明 $ (A \eiff A) \eand (B \eand \enot B)$ 是矛盾式。
\item 证明 $ (B \eiff (C \eor B)) $ 不是矛盾式。
\item 证明 $ ((\enot X \eiff X) \eor X) $ 与 $X$ 逻辑等价。
\item 证明 $F \eand (K \eand R) $ 与 $ (F \eiff (K \eiff R)) $ 不逻辑等价。
\item 证明语句 $ \enot (W \eif W)$、$(W \eiff W) \eand W$、$E \eor (W \eif \enot (E \eand W))$ 是共同不可满足的。
\item 证明语句  $\enot R \eor C $、$(C \eand R) \eif \enot R$、$(\enot (R \eor C) \eif R) $ 是共同可满足的。
\item 证明 $\enot \enot (C \eiff \enot C), ((G \eor C) \eor G) \therefore ((G \eif C) \eand G) $ 是有效的。
\item 证明 $ \enot \enot L,  (C \eif \enot L) \eif C) \therefore \enot C$ 是无效的。
\end{compactlist}
